% =====================================================================
%  Informe tecnico consolidado - Instancia ATENEA2 (SRVPRORPDB01)
%  Departamento de IT - Servicios Generales de la Plana
%  Compilar con:  pdflatex informe-atenea2-consolidado.tex   (dos pasadas)
% =====================================================================
\documentclass[11pt,a4paper]{article}

\usepackage[utf8]{inputenc}
\usepackage[T1]{fontenc}
\usepackage[spanish,es-noquoting,es-nolists]{babel}
\usepackage[a4paper,margin=2.3cm]{geometry}
\usepackage{lmodern}
\usepackage{booktabs}
\usepackage{longtable}
\usepackage{array}
\usepackage[table]{xcolor}
\usepackage{listings}
\usepackage{graphicx}
\usepackage{tikz}
\usepackage{pgfplots}
\pgfplotsset{compat=1.18}
\usepackage[hidelinks]{hyperref}
\usepackage{fancyhdr}
\usepackage{titlesec}
\usepackage{enumitem}

% ----------------------------- colores -------------------------------
\definecolor{tinta}{HTML}{1A1A1A}
\definecolor{acento}{HTML}{1F6FEB}
\definecolor{gris}{HTML}{6B7280}
\definecolor{grisclaro}{HTML}{F3F4F6}
\definecolor{rojo}{HTML}{B91C1C}
\definecolor{verde}{HTML}{15803D}
\definecolor{ambar}{HTML}{B45309}
\definecolor{comentario}{HTML}{6B7280}
\definecolor{cadena}{HTML}{15803D}

% --------------------------- cabecera/pie ----------------------------
\pagestyle{fancy}
\setlength{\headheight}{14pt}
\setlength{\tabcolsep}{4pt}
\fancyhf{}
\renewcommand{\headrulewidth}{0.4pt}
\renewcommand{\footrulewidth}{0pt}
\fancyhead[L]{\footnotesize\color{gris} ATENEA2 · SRVPRORPDB01 · Diagnóstico de bloqueos}
\fancyhead[R]{\footnotesize\color{gris}\thepage}
\fancyfoot[C]{\footnotesize\color{gris} SGP · Departamento de IT}

% --------------------------- titulos ---------------------------------
\titleformat{\section}{\Large\bfseries\color{tinta}}{\thesection}{0.6em}{}
\titleformat{\subsection}{\large\bfseries\color{tinta}}{\thesubsection}{0.6em}{}
\titleformat{\subsubsection}{\normalsize\bfseries\color{gris}}{}{0em}{}

% --------------------------- listings --------------------------------
\lstdefinestyle{sql}{
  language=SQL, basicstyle=\ttfamily\scriptsize, keywordstyle=\color{acento}\bfseries,
  commentstyle=\color{comentario}\itshape, stringstyle=\color{cadena},
  numbers=left, numberstyle=\tiny\color{gris}, numbersep=8pt,
  backgroundcolor=\color{grisclaro}, frame=leftline, rulecolor=\color{gris},
  framexleftmargin=6pt, xleftmargin=14pt, breaklines=true,
  showstringspaces=false, columns=fullflexible, morekeywords={APPLY,CROSS,OVER,PARTITION,WITH,STUFF,DATEADD,DATEDIFF,CAST,NULLIF,ISNULL,TOP}
}
\lstdefinestyle{salida}{
  basicstyle=\ttfamily\scriptsize\color{tinta}, backgroundcolor=\color{grisclaro},
  frame=leftline, rulecolor=\color{gris}, framexleftmargin=6pt,
  xleftmargin=10pt, breaklines=true, showstringspaces=false, columns=fullflexible
}

\newcommand{\cod}[1]{\texttt{\small #1}}
\newcommand{\critico}[1]{\textcolor{rojo}{\textbf{#1}}}
\newcommand{\bien}[1]{\textcolor{verde}{\textbf{#1}}}
\newcommand{\aviso}[1]{\textcolor{ambar}{\textbf{#1}}}
\newcommand{\sombra}{\cellcolor{grisclaro}}

% =====================================================================
\begin{document}

% --------------------------- PORTADA ---------------------------------
\begin{titlepage}
\centering
\vspace*{2.5cm}

{\color{gris}\large SERVICIOS GENERALES DE LA PLANA · DEPARTAMENTO DE IT}

\vspace{1.4cm}
{\Huge\bfseries\color{tinta} Diagnóstico de bloqueos\\[0.35cm] en la instancia ATENEA2}

\vspace{0.8cm}
{\large\color{gris} SRVPRORPDB01 · SQL Server 2022 · 16 vCPU · 128\,GB}

\vspace{2.2cm}
\begin{minipage}{0.82\textwidth}
\color{tinta}\normalsize
Informe técnico consolidado. Reúne todas las mediciones realizadas sobre la
instancia entre el 22 y el 29 de octubre de 2026, los hallazgos confirmados,
las hipótesis descartadas y el plan de corrección por fases.

\vspace{0.5cm}
Incluye además el análisis de la base de datos \cod{Cenusa}, que apareció
durante la investigación y resultó concentrar los costes unitarios más altos
de toda la instancia.
\end{minipage}

\vfill
\begin{minipage}{0.82\textwidth}
\color{gris}\small
\begin{tabular}{@{}ll@{}}
\textbf{Instancia}   & SRVPRORPDB01 (ATENEA2), SQL Server 2022, compat. 160 \\
\textbf{Fuentes}     & SolarWinds DPA 2022.4 · consultas directas · vCenter \\
\textbf{Cuenta}      & \cod{dpaMONITOR} (solo lectura) \\
\textbf{Ventana}     & 22--29 de octubre de 2026 \\
\textbf{Versión}     & 4.0 — consolidada \\
\textbf{Fecha}       & 29 de octubre de 2026 \\
\end{tabular}
\end{minipage}

\vspace{1cm}
\end{titlepage}

% --------------------- HISTORIAL DE REVISIONES -----------------------
\section*{Historial de revisiones}
\addcontentsline{toc}{section}{Historial de revisiones}

La tesis central de este informe cambió tres veces durante la investigación.
Queda registrado porque afecta a cómo deben leerse las recomendaciones.

\begin{longtable}{p{1.1cm}p{2.1cm}p{12.4cm}}
\toprule
\textbf{Ver.} & \textbf{Fecha} & \textbf{Cambio} \\
\midrule
\endhead
1.0 & 29/10/2026 & Primera versión. Tesis: \emph{paralelismo mal
  configurado} (se suponía MAXDOP 0 y \emph{cost threshold} 5). \\
2.0 & 29/10/2026 & \critico{Tesis 1.0 refutada} por medición directa: MAXDOP
  estaba en 8 y el umbral en 100. Nueva tesis: paralelismo \emph{innecesario}
  por estimaciones de cardinalidad erróneas. Se añade el análisis del SQL
  Agent. \\
3.0 & 29/10/2026 & Se registra el vCenter en DPA y se descarta la contención
  de hipervisor. Se añaden latencia por fichero e índices ausentes. \\
4.0 & 29/10/2026 & Consolidación. Se demuestra la tesis 2.0 con evidencia de
  planes. Se añaden escalado de bloqueos, grupo de disponibilidad (descartado)
  y la base de datos \cod{Cenusa}. \critico{Se corrige} la lectura de
  \cod{HADR\_NOTIFICATION\_DEQUEUE} y \critico{se descarta} la hipótesis de
  \cod{IndexOptimize} como causa del bloqueo de esquema. \\
\bottomrule
\end{longtable}

\newpage
\tableofcontents
\newpage

% =====================================================================
\section{Resumen ejecutivo}

\subsection*{Qué se investigó}

La instancia ATENEA2 presenta sesiones bloqueadas de forma \textbf{permanente
durante toda la jornada laboral}. En los siete días analizados, el
\critico{88\,\%} de las muestras registran bloqueos, con una media de 4,9
sesiones bloqueadas y un pico de 9,0 a las 08:00. Las otras cuatro instancias
monitorizadas del parque se mueven entre el 0 y el 2\,\%.

\subsection*{Conclusión en cuatro frases}

\begin{enumerate}[leftmargin=1.4em]\itemsep3pt
  \item \textbf{La configuración de SQL Server es correcta} y la plataforma
        también: paralelismo, TempDB, disco e hipervisor se midieron y se
        descartaron uno por uno.
  \item \textbf{El problema es de estimación, no de configuración}: el
        optimizador trabaja con estadísticas permanentemente obsoletas, elige
        planes paralelos que no hacen falta y concede memoria que no se usa.
  \item \textbf{La concurrencia lo amplifica}: 140 planificaciones de SQL
        Agent sub-diarias, ninguna base de datos con versionado de filas, y
        575 disparadores en \cod{Cenusa} que alargan cada transacción desde
        dentro.
  \item \textbf{No hay una sola causa}, y por eso no hay un solo arreglo: las
        catorce acciones van en tres fases, de lo reversible y gratuito a lo
        que exige ventana y desarrollo.
\end{enumerate}

\subsection*{Los seis hallazgos confirmados}

\begin{longtable}{p{0.4cm}p{6.6cm}p{8.6cm}}
\toprule
& \textbf{Hallazgo} & \textbf{Evidencia} \\
\midrule
\endhead
1 & Estadísticas permanentemente obsoletas
  & \cod{entEFECTO\_FacturasCobradas365}: 31,5\,M de modificaciones sobre
    1,18\,M de filas en 9\,h (26$\times$ su tamaño), muestreada al 11,8\,\% \\
2 & Paralelismo innecesario, demostrado
  & 6 variantes del mismo \cod{INSERT}: 31.000\,s de CPU para devolver
    \textbf{0 filas}. Una consulta de \cod{OMS}: 506\,MB concedidos para
    devolver \textbf{1 fila} \\
3 & Escalado y espera de bloqueos concentrados
  & \cod{entEMPRESA}: \textbf{38 horas} de espera acumulada en 27 días.
    \cod{entORDENCARGA}: 22 horas. Tres jobs reescriben \cod{entEMPRESA}
    cada 30 minutos en horario laboral \\
4 & Sin versionado de filas
  & RCSI desactivado en las 8 bases revisadas: los lectores bloquean a los
    escritores \\
5 & SQL Agent saturado y desbordando la noche
  & 140 planificaciones sub-diarias; 2 jobs al 118\,\% y 96\,\% de ocupación;
    \cod{TAREAS NOCTURNAS STN} puede acabar a las 09:14 \\
6 & \cod{Cenusa}: 575 disparadores
  & 9 sobre \cod{entEMPRESA}. \cod{AUDITORIA2}: 29,7\,M de filas con 89\,M de
    modificaciones, muestreada al 1\,\%. Un procedimiento consume
    \textbf{1,1\,GB} de memoria y aprovecha el 17\,\% \\
\bottomrule
\end{longtable}

\subsection*{Los trece descartes}

Medidos y eliminados: configuración de paralelismo (a nivel de instancia y de
base de datos), estimador de cardinalidad, dimensionamiento de TempDB,
almacenamiento (desde el invitado y desde el host), proceso ajeno consumiendo
CPU, agotamiento del grupo de hilos, contención de CPU del hipervisor,
presión de memoria del hipervisor, confirmación síncrona del grupo de
disponibilidad, índices ausentes en \cod{Cenusa} y \cod{IndexOptimize} como
causa del bloqueo de esquema. \textbf{Ninguno se descartó por criterio: todos
por medición directa.} El detalle, con la consulta y el resultado de cada
uno, en la sección \ref{sec:descartes}.

\subsection*{Qué contiene este informe}

\begin{longtable}{p{1.1cm}p{4.6cm}p{9.9cm}}
\toprule
\textbf{Sec.} & \textbf{Título} & \textbf{Contenido} \\
\midrule
\endhead
2 & Alcance y método & De dónde sale cada dato, con qué ventana temporal, y
    las cuatro correcciones sobre el propio análisis \\
3 & El síntoma & Qué se observó: capturas del incidente, composición de las
    sesiones, recurrencia de 7 días y contraste con el resto del parque \\
4 & Diagnóstico & El mecanismo completo: diagrama causal, esperas, la prueba
    del paralelismo, estadísticas, escalado de bloqueos, RCSI, SQL Agent,
    memoria, CPU, hipervisor y registros de transacciones \\
5 & \cod{Cenusa} & Los 575 disparadores, \cod{AUDITORIA2}, el procedimiento
    de 1,1\,GB y qué está bien en esa base \\
6 & Descartes & Tabla consolidada: hipótesis, qué se midió, resultado \\
7 & Plan de acción & Doce acciones en tres fases, cada una con su métrica de
    verificación. Y lo que no se recomienda hacer \\
8 & Consultas & Todo el SQL utilizado, reproducible tal cual \\
9 & Evidencia & Salida literal de cada consulta, sin editar \\
10 & Anexo & Las 80 planificaciones de menos de una hora \\
11 & Glosario & Términos técnicos empleados \\
\bottomrule
\end{longtable}

\vspace{0.4cm}
\noindent\fbox{\parbox{0.975\textwidth}{\textbf{Lo que yo haría el lunes:}
activar \cod{AUTO\_UPDATE\_STATISTICS\_ASYNC} en \cod{Grupo} y programar
estadísticas con \cod{FULLSCAN} tres o cuatro veces al día sobre
\cod{entEFECTO\_FacturasCobradas365} y
\cod{STN\_ERP\_ANALISIS\_VENTAS\_PEDIDOS}. Es un comando y un job, sin
ventana de parada, reversible en segundos, y ataca el mecanismo que genera
todo lo demás. \textbf{Un cambio cada 48 horas}, nunca varios a la vez.}}

% =====================================================================
\newpage
\section{Alcance y método}

\subsection{De dónde sale cada dato}

Tres fuentes, con capacidades y ventanas muy distintas. Distinguirlas importa
porque no todas las cifras de este informe tienen la misma solidez.

\begin{longtable}{p{0.4cm}p{4.3cm}p{5.1cm}p{5.4cm}}
\toprule
& \textbf{Fuente} & \textbf{Qué aporta} & \textbf{Ventana y granularidad} \\
\midrule
\endhead
1 & Panel propio en
    \cod{http://10.50.69.121:9000}, que lee de SolarWinds DPA 2022.4
  & Histórico de sesiones activas y bloqueadas, esperas, CPU, memoria,
    TempDB, backups
  & 22--29/10. Diezmado a 720 puntos por serie: 1 muestra/hora en la ventana
    de 7 días \\
2 & Consultas directas a SRVPRORPDB01 por la herramienta de consultas de DPA,
    con \cod{dpaMONITOR}
  & Vistas de gestión dinámica: configuración, esperas, bloqueos,
    estadísticas, índices, jobs, planes
  & Acumulado desde el arranque (02/10), 27 días. Solo lectura \\
3 & Pestaña \emph{Virtualization} de DPA, alimentada por el vCenter
    \cod{vcsa.central.local}
  & \emph{Ready time}, \emph{co-stop}, \emph{ballooning}, intercambio del
    host, latencia por dispositivo
  & \critico{Solo $\approx$70 minutos}: el vCenter se registró el mismo día
    29/10 \\
\bottomrule
\end{longtable}

\subsection{La instancia}

\begin{table}[h]
\centering
\small
\begin{tabular}{ll}
\toprule
\textbf{Característica} & \textbf{Valor} \\
\midrule
Servidor              & SRVPRORPDB01 (ATENEA2), VM sobre ESXi \cod{10.150.0.8} \\
Versión               & SQL Server 2022, nivel de compatibilidad 160 \\
Núcleos               & 16 vCPU en 2 nodos NUMA, 16 \emph{schedulers} \\
Memoria               & 131.071\,MB físicos; \cod{max server memory} 96.896\,MB \\
Arranque              & 02/10/2026 14:25 — 27 días de actividad acumulada \\
Alta disponibilidad   & \cod{AG-ATENEA2}, réplica secundaria SRVPRORPDB02 \\
Bases relevantes      & \cod{Grupo}, \cod{Cenusa}, \cod{OMS}, \cod{IFR},
                        \cod{ContaSTN}, \cod{D365Eventos} \\
\bottomrule
\end{tabular}
\end{table}

\subsection{Limitaciones conocidas}

\paragraph{El histórico del panel se diezma.} Cada serie se limita a 720
puntos, así que la ventana de siete días se reduce a una muestra por hora y
la de 24 horas a una cada diez minutos. \textbf{Los máximos citados son cotas
inferiores}: los picos entre muestras no se ven.

\paragraph{El contador de bloqueos y su lista no son simultáneos.} El
contador procede de una métrica que DPA muestrea cada minuto; la lista de
sesiones es una foto en vivo tomada en otro instante. Que el contador diga 23
y la lista muestre 3 no es un error.

\paragraph{Las vistas acumuladas no distinguen la hora.}
\cod{dm\_os\_wait\_stats}, \cod{dm\_io\_virtual\_file\_stats},
\cod{dm\_db\_index\_operational\_stats} y \cod{dm\_db\_missing\_index\_*}
suman los 27 días desde el arranque sin separar las horas punta de la
madrugada. Sirven para saber \emph{qué} domina, no \emph{cuándo} ocurre. La
datación de este informe viene del panel, no de ellas.

\paragraph{Los índices sugeridos son heurísticos.} Las cifras de impacto de
\cod{dm\_db\_missing\_index\_group\_stats} las estima el optimizador sobre
planes individuales, ignorando los índices existentes y el coste de
mantenimiento en escritura. Indican dónde mirar; no son un guion para
aplicar.

\paragraph{El detalle de sesiones se refresca cada dos minutos} y se limita a
60 filas ordenadas por tiempo de espera. Los repartos por programa, host y
base de datos salen de esa muestra, no del censo completo.

\paragraph{El hipervisor tiene 70 minutos de histórico.} Las magnitudes son
inequívocas, pero \textbf{el pico de las 08:00 no se ha observado todavía a
ese nivel}. Procede repetir la comprobación con 7 días completos.

\paragraph{Query Store está desactivado en \cod{Grupo}} (solo activo en
\cod{D365Eventos}). Sin él no hay histórico de planes, que es lo que
permitiría seguir la evolución de las estimaciones en el tiempo en lugar de
verlas en un corte.

\paragraph{Queda sin identificar el origen del bloqueo \cod{Sch-M} de
392\,s.} Ver la sección \ref{sec:schm}.

\paragraph{\cod{Cenusa} no se ha agotado.} Se analizaron sus disparadores,
estadísticas, índices ausentes y esperas de bloqueo, pero \textbf{no el
contenido funcional de los 575 disparadores}, que es trabajo de desarrollo.

\subsection{Cuatro correcciones sobre el propio análisis}
\label{sec:correcciones}

Durante la investigación se descartaron cuatro afirmaciones propias. Se
registran porque definen cuánto crédito merece el razonamiento frente a la
medición.

\begin{longtable}{p{0.4cm}p{6.4cm}p{8.8cm}}
\toprule
& \textbf{Se afirmó} & \textbf{Y resultó} \\
\midrule
\endhead
1 & El paralelismo estaba mal configurado: MAXDOP 0 y \emph{cost threshold}
    en el valor por defecto de 5
  & \critico{Falso.} MAXDOP 8 y umbral 100, ambos correctos. También se
    comprobó el ámbito de base de datos, sin sobrescritura. La tesis se
    sustituyó por «paralelismo innecesario», con un remedio completamente
    distinto: estadísticas, no configuración \\
2 & El log de TempDB estaba al 0\,\% libre, según el panel
  & \critico{Falso.} Medido directamente, al 93\,\% libre. Falsa alarma de esa
    métrica de DPA \\
3 & \cod{HADR\_NOTIFICATION\_DEQUEUE} (2,8\,\% de las esperas) indicaba carga
    del grupo de disponibilidad
  & \critico{Falso.} Es una espera \emph{inactiva} de hilos de fondo. La
    sospecha derivada —AG en confirmación síncrona— se comprobó y también
    resultó falsa: es \cod{ASYNCHRONOUS\_COMMIT} \\
4 & El bloqueo \cod{Sch-M} de 392\,s lo causaba \cod{IndexOptimize}
  & \critico{Falso.} Ese día el job terminó a las 03:26:58, siete horas antes
    del incidente \\
\bottomrule
\end{longtable}

\noindent Cada acción de la sección \ref{sec:plan} lleva por ello su propia
consulta de verificación. \textbf{Las recomendaciones valen lo que valen las
mediciones que las sostienen}, no la plausibilidad del razonamiento.

% =====================================================================
\newpage
\section{El síntoma}

\subsection{Qué se observó}

Dos capturas del mismo día, separadas por minutos, con el panel como fuente.

\begin{table}[h]
\centering
\small
\begin{tabular}{lrr}
\toprule
\textbf{Medida} & \textbf{Captura 1} & \textbf{Captura 2} \\
\midrule
Sesiones activas              & 202 & 187 \\
Sesiones bloqueadas (contador)& \critico{23} & \critico{18} \\
Espera acumulada en 1 hora    & \critico{1.527\,s} & --- \\
Cadena de bloqueo más larga   & 3 niveles & 2 niveles \\
Espera máxima de una sesión   & \critico{392\,s} & 214\,s \\
\bottomrule
\end{tabular}
\end{table}

Una sesión esperó \textbf{392 segundos} —seis minutos y medio— por un bloqueo
de esquema. Se observaron esperas de 200 segundos en clientes
\cod{.Net SqlClient} desde puestos de usuario. Eso no se percibe como
lentitud: se percibe como que la aplicación se ha colgado, y el usuario
cierra y vuelve a entrar, añadiendo otra sesión a la cola.

\subsection{Composición de las sesiones}

De las 60 sesiones examinadas en detalle (la muestra que devuelve el panel,
ordenada por tiempo de espera):

\begin{table}[h]
\centering
\small
\begin{tabular}{lrl}
\toprule
\textbf{Esperando por} & \textbf{Sesiones} & \textbf{Interpretación} \\
\midrule
Paralelismo (\cod{CX*}, \cod{EXECSYNC}) & \critico{45} & reparto de trabajo entre hilos \\
Bloqueo de objetos (\cod{LCK\_M\_*})    & 7 & espera real por bloqueo \\
CPU (\cod{SOS\_SCHEDULER\_YIELD})       & 5 & turno de procesador \\
Otros                                   & 3 & --- \\
\bottomrule
\end{tabular}
\end{table}

\textbf{El dato clave está en el reparto.} Solo 7 de 60 esperaban por un
bloqueo; 45 esperaban por paralelismo. Y las sesiones que bloqueaban a las
demás no esperaban a nadie: esperaban turno de CPU. Es decir, \textbf{el
bloqueo no es el origen, es la consecuencia} de que los bloqueadores tarden
demasiado en soltar.

Se detectaron además \textbf{10 sesiones dormidas} (\cod{sleeping}) con
transacción abierta, reteniendo bloqueos sin ejecutar nada: indicio de
transacciones que la aplicación abre y no cierra con diligencia.

\subsection{Qué tan recurrente es}

\begin{figure}[h]
\centering
\begin{tikzpicture}
\begin{axis}[
  width=15.2cm, height=5.6cm,
  xlabel={\footnotesize Hora del día}, ylabel={\footnotesize Sesiones bloqueadas},
  xmin=0, xmax=23, ymin=0, ymax=10,
  xtick={0,2,4,6,8,10,12,14,16,18,20,22},
  ytick={0,2,4,6,8,10},
  grid=major, grid style={gris!25},
  tick label style={font=\scriptsize\color{gris}},
  label style={font=\scriptsize\color{gris}},
  axis lines=left, every axis plot/.append style={thick},
  legend style={font=\scriptsize, draw=none, fill=none, at={(0.98,0.97)}, anchor=north east}
]
\addplot[color=acento, mark=*, mark size=1.3pt] coordinates {
 (0,2.1)(1,2.4)(2,3.1)(3,2.8)(4,3.4)(5,4.2)(6,5.8)(7,7.4)(8,9.0)(9,8.1)
 (10,6.9)(11,6.2)(12,5.4)(13,4.1)(14,4.6)(15,5.9)(16,6.1)(17,5.2)(18,4.4)
 (19,3.6)(20,3.1)(21,2.7)(22,2.3)(23,2.0)};
\addlegendentry{media por hora (7 días)}
\addplot[color=rojo, dashed, no marks] coordinates {(0,1)(23,1)};
\addlegendentry{objetivo: menos de 1}
\node[font=\scriptsize\color{rojo}, anchor=south] at (axis cs:8,9.0) {9,0};
\end{axis}
\end{tikzpicture}
\caption{Media de sesiones bloqueadas por hora, del 22 al 29 de octubre.
El pico de las 08:00 coincide con la entrada de los usuarios y con el final
de los procesos nocturnos que se desbordan.}
\end{figure}

\begin{table}[h]
\centering
\small
\begin{tabular}{lr}
\toprule
\textbf{Indicador (7 días)} & \textbf{Valor} \\
\midrule
Muestras con al menos una sesión bloqueada & \critico{88\,\%} \\
Media de sesiones bloqueadas               & 4,9 \\
Pico horario (08:00)                       & \critico{9,0} \\
\emph{Page life expectancy} medio          & \critico{124\,s} \\
\emph{Page life expectancy} instantáneo    & \critico{28\,s} \\
\bottomrule
\end{tabular}
\end{table}

\textbf{No es un incidente puntual.} El 88\,\% de las muestras registran
bloqueo: es el estado normal de esta instancia durante la jornada.

\subsection{Contraste con el resto del parque}

\begin{table}[h]
\centering
\small
\begin{tabular}{llrr}
\toprule
\textbf{Instancia} & \textbf{Versión} & \textbf{\% muestras con bloqueo} &
\textbf{Media} \\
\midrule
\critico{ATENEA2} (SRVPRORPDB01) & SQL 2022 & \critico{88\,\%} & \critico{4,9} \\
SCALEPROD (SRVPROSCDB1)          & SQL 2022 & 2\,\% & 0,1 \\
SRVPBI                           & SQL 2022 & 1\,\% & 0,0 \\
PROMETEO                         & SQL 2014 & 0\,\% & 0,0 \\
AQUILES                          & SQL 2012 & 0\,\% & 0,0 \\
\bottomrule
\end{tabular}
\end{table}

Las cuatro restantes son irrelevantes estadísticamente. \textbf{El problema
es específico de ATENEA2}, no una característica del parque ni de la versión
de SQL Server: SCALEPROD corre la misma versión sobre hardware comparable y
está al 2\,\%.

% =====================================================================
\newpage
\section{Diagnóstico: el mecanismo}

\subsection{El cuadro completo}

No hay una causa única. Hay \textbf{tres condiciones de fondo} que se
combinan, y un ciclo que una vez arrancado se sostiene solo.

\begin{figure}[h]
\centering
\begin{tikzpicture}[
  caja/.style={draw=gray!60, rounded corners=4pt, minimum width=4.4cm,
               minimum height=1.35cm, align=center, font=\footnotesize},
  raiz/.style={caja, draw=rojo!70, line width=1pt, fill=rojo!6,
               minimum height=1.6cm},
  destacada/.style={caja, draw=acento, line width=1.2pt, fill=acento!8},
  flecha/.style={-latex, draw=gray!70, line width=0.9pt},
  fina/.style={-latex, draw=rojo!50, line width=0.8pt}
]
\node[raiz] (est)  at (-5.2, 4.8) {\textbf{Estadísticas obsoletas}\\
    \scriptsize\color{gris} 31,5\,M de cambios sobre\\[-1pt]
    \scriptsize\color{gris} 1,18\,M de filas};
\node[raiz] (job)  at ( 0.0, 4.8) {\textbf{140 jobs sub-diarios}\\
    \scriptsize\color{gris} dos al 118\,\% y al 96\,\%\\[-1pt]
    \scriptsize\color{gris} de ocupación};
\node[raiz] (rcsi) at ( 5.2, 4.8) {\textbf{Sin versionado de filas}\\
    \scriptsize\color{gris} RCSI desactivado\\[-1pt]
    \scriptsize\color{gris} en las ocho bases};

\node[destacada] (plan) at (-3.6, 2.2) {\textbf{Estimaciones erróneas}\\
    \scriptsize\color{gris} planes paralelos innecesarios};
\node[caja]      (escr) at ( 3.6, 2.2) {\textbf{Escritura constante}\\
    \scriptsize\color{gris} y lectores que bloquean};

\node[destacada] (cpu)  at (-3.6, -0.3) {\textbf{CPU saturada}\\
    \scriptsize\color{gris} 72--83\,\%; 58,7\,\% de esperas\\[-1pt]
    \scriptsize\color{gris} son de reparto paralelo};
\node[caja]      (tran) at ( 3.6, -0.3) {\textbf{Transacciones largas}\\
    \scriptsize\color{gris} retienen bloqueos más tiempo};

\node[caja]      (bloq) at ( 0.0, -2.8) {\textbf{Bloqueo permanente}\\
    \scriptsize\color{gris} 88\,\% de las muestras, pico de 9 a las 8\,h};

\draw[fina]   (est)  -- (plan);
\draw[fina]   (job)  -- (escr);
\draw[fina]   (rcsi) -- (escr);
\draw[flecha] (plan) -- (cpu);
\draw[flecha] (escr) -- (tran);
\draw[flecha] (cpu)  -- (tran);
\draw[flecha] (cpu)  -- (bloq);
\draw[flecha] (tran) -- (bloq);
\draw[flecha] (bloq) to[out=200,in=250,looseness=1.6]
      node[below left,font=\scriptsize,color=gris,pos=0.45]{se realimenta} (cpu);
\end{tikzpicture}
\caption{Tres condiciones de fondo (en rojo) alimentan un ciclo que se
sostiene a sí mismo. El hipervisor, el disco y la configuración de
paralelismo quedaron descartados por medición.}
\end{figure}

La realimentación explica por qué el cuadro no se corrige solo: las sesiones
que esperan turno de CPU \textbf{retienen sus bloqueos mientras esperan}, y
los usuarios que ven la aplicación colgada cierran y vuelven a entrar,
añadiendo sesiones a la misma cola.

\subsection{Las esperas}

Acumulado de los 27 días que lleva encendida la instancia:

\begin{table}[h]
\centering
\small
\begin{tabular}{lrrl}
\toprule
\textbf{Tipo de espera} & \textbf{Segundos} & \textbf{\% total} &
\textbf{Qué indica} \\
\midrule
CXCONSUMER                  & 45.504.691 & 27,6\,\% & paralelismo \\
CXSYNC\_PORT                & 31.022.696 & 18,8\,\% & paralelismo \\
CXPACKET                    & 20.425.851 & 12,4\,\% & paralelismo \\
SOS\_SCHEDULER\_YIELD       & 12.906.011 &  7,8\,\% & \critico{falta de CPU} \\
LATCH\_EX                   &  9.896.832 &  6,0\,\% & contención interna \\
\sombra HADR\_NOTIFICATION\_DEQUEUE & \sombra 4.645.761 & \sombra 2,8\,\% &
  \sombra \emph{inactiva --- no computa} \\
HTDELETE                    & --- &  2,2\,\% & uniones hash en paralelo \\
OLEDB                       & --- &  2,1\,\% & consultas distribuidas \\
PAGEIOLATCH\_SH             & --- &  2,1\,\% & \critico{lecturas a disco por falta de memoria} \\
\bottomrule
\end{tabular}
\end{table}

\paragraph{Corrección (ver \ref{sec:correcciones}).}
\cod{HADR\_NOTIFICATION\_DEQUEUE}, como \cod{HADR\_WORK\_QUEUE},
\cod{HADR\_TIMER\_TASK} y \cod{HADR\_FILESTREAM\_IOMGR\_IOCOMPLETION},
corresponde a hilos de fondo ociosos esperando trabajo. Aparece con cifras
enormes en cualquier instancia con grupo de disponibilidad y debe excluirse,
igual que \cod{SLEEP\_*} o \cod{XE\_TIMER\_EVENT}. En una versión anterior se
interpretó como carga del AG: era un error.

\paragraph{El paralelismo suma el \critico{58,7\,\%} de todo el tiempo de
espera}, pese a estar bien configurado. Durante buena parte de la
investigación se asumió que la causa era «demasiado trabajo pesado a la vez
para 16 núcleos». La sección siguiente demuestra que no: \textbf{se reparte
trabajo que no hacía falta repartir}.

\subsection{La prueba: planes paralelos que no devuelven nada}
\label{sec:prueba}

La huella de una estimación inflada es doble: el optimizador \textbf{concede
memoria} para las filas que cree que va a manejar, y \textbf{reparte la
consulta entre varios núcleos} por el mismo motivo. Si después devuelve cero
filas, la estimación era errónea. Consultando
\cod{sys.dm\_exec\_query\_stats}:

\begin{table}[h]
\centering
\small
\begin{tabular}{lrrrrr}
\toprule
\textbf{Sentencia (BD \cod{Grupo})} & \textbf{Ejec.} & \textbf{CPU} &
\textbf{Duración} & \textbf{Núcleos} & \textbf{Filas} \\
\midrule
\cod{INSERT StnERPArticuloLoteLinea} ($\times$6) & 2.147 & \critico{31.000\,s} &
  12.166\,s & 2,5 & \critico{0} \\
\cod{UPDATE stnERPORDENPREPARACION} & 1.122 & 1.340\,s & 376\,s &
  \critico{3,6} & \critico{0} \\
\cod{UPDATE stnERPORDENPREPARACION} (2) & 1.122 & 919\,s & 328\,s & 2,8 &
  \critico{0} \\
\cod{DELETE D365Eventos..STNAppSuite\ldots} & 591 & 1.993\,s & 898\,s & 2,2 & 3 \\
\cod{DELETE D365Eventos..STNAppSuite\ldots} (2) & 597 & 1.255\,s & 597\,s & 2,1 & 3 \\
\cod{DELETE ..InventTransOriginBiEntities} & 595 & 1.181\,s & 627\,s & 1,9 & 2 \\
\cod{SELECT OMS..InventTransOrigin\ldots} & 164 & 581\,s & 291\,s & 2,0 &
  \critico{1} \\
\bottomrule
\end{tabular}
\caption{«Núcleos» = \cod{total\_worker\_time / total\_elapsed\_time}: por
encima de 1 indica plan paralelo. «Filas» es la media devuelta por ejecución.}
\end{table}

\paragraph{Ocho horas y media de CPU para no devolver ninguna fila.} Las seis
variantes de \cod{INSERT INTO StnERPArticuloLoteLinea} suman \textbf{31.000
segundos de CPU} en apenas 2.147 ejecuciones: 14 segundos de CPU cada una,
repartidos entre 2,5 núcleos, con 71\,MB de memoria concedida y cero filas de
resultado. Es el mayor consumidor paralelo de la instancia.

\paragraph{506\,MB de memoria para devolver una fila.}
\cod{SELECT a.SourceKey FROM OMS..InventTransOriginBiEntities} recibe
\textbf{506\,MB de concesión por ejecución} y devuelve \textbf{una fila}. No
hay forma más clara de enseñar una estimación de cardinalidad equivocada. Esa
concesión se resta del \emph{buffer pool} mientras dura la consulta, lo que
conecta directamente con el \emph{page life expectancy} de 28 segundos.

\paragraph{Seis copias de la misma sentencia.} Que el mismo \cod{INSERT}
aparezca seis veces con estadísticas propias indica que \textbf{no está
parametrizado}: cada variante literal genera su propio plan, ocupa su propia
entrada en la caché y se optimiza por separado.

\paragraph{Conclusión.} El \emph{cost threshold} está en 100 y es correcto,
pero es un umbral sobre un coste \emph{estimado}. Si la estimación de filas
está inflada, se cruza igual. \textbf{El optimizador no está mal configurado:
está eligiendo paralelismo para consultas que no lo necesitan.}

\subsection{Las estadísticas están permanentemente obsoletas}

\begin{table}[h]
\centering
\small
\begin{tabular}{lrrrr}
\toprule
\textbf{Tabla (BD \cod{Grupo})} & \textbf{Filas} & \textbf{Modificaciones} &
\textbf{Churn} & \textbf{Muestra} \\
\midrule
entEFECTO\_FacturasCobradas365 & 1.184.877 & \critico{31.476.588} &
  \critico{26,6$\times$} & \critico{11,8\,\%} \\
STN\_ERP\_ANALISIS\_VENTAS\_PEDIDOS & 6.460.673 & \critico{12.940.889} &
  \critico{2,0$\times$} & \critico{1,4\,\%} \\
\bottomrule
\end{tabular}
\caption{Estadísticas actualizadas esa misma madrugada, a las 02:41 y 02:48.}
\end{table}

Las estadísticas se actualizaron \emph{hoy}. No son antiguas por fecha
(\cod{dias} = 0). El problema es el volumen de cambio acumulado desde
entonces: en menos de nueve horas,
\cod{entEFECTO\_FacturasCobradas365} registró \textbf{31,5 millones de
modificaciones sobre una tabla de 1,18 millones de filas}. Veintiséis veces
su propio tamaño. Es una tabla de paso que se vacía y recarga continuamente,
lo que concuerda con los jobs \cod{D365 CARGAR FACTURAS COBRADAS} (cada hora),
\cod{ACTUALIZAR\_FECHACOBRO\_EFECTOS\_DESDE\_CONCILIADOS} (cada 30 minutos) y
\cod{Limpiar Fechas Efectos} (cada hora).

El umbral de actualización automática en nivel 160 es aproximadamente
$\sqrt{1000 \times filas}$, unas 34.400 modificaciones para esta tabla. Se
cruza en minutos. En la práctica, \textbf{el optimizador nunca trabaja con
estadísticas representativas de esa tabla}.

Y el muestreo agrava el efecto: 11,8\,\% en una, \textbf{1,4\,\%} en otra de
6,4 millones de filas. Estimar la cardinalidad de 6,4 millones de filas a
partir del 1,4\,\%, con el doble de su tamaño en cambios pendientes, produce
estimaciones muy alejadas de la realidad. \textbf{Esta es la causa raíz de lo
demostrado en \ref{sec:prueba}}.

\paragraph{Agravante: las actualizaciones son sincrónicas.}
\cod{is\_auto\_update\_stats\_async\_on} está en \cod{false} en las ocho
bases revisadas. Cuando se cruza el umbral, la consulta que lo cruza
\textbf{espera} a que se recalculen las estadísticas antes de compilar. Sobre
una tabla de 6,4 millones de filas, en horario laboral, con 200 sesiones
activas, esa espera se la come el usuario.

\subsection{Escalado de bloqueos: dónde duele de verdad}
\label{sec:escalado}

\cod{sys.dm\_db\_index\_operational\_stats} mide, por tabla, cuántas veces
SQL Server convirtió bloqueos de fila en bloqueos de tabla y cuánto tiempo se
esperó por bloqueos. Acumulado de 27 días en \cod{Grupo}:

\begin{table}[h]
\centering
\small
\begin{tabular}{lrrrr}
\toprule
\textbf{Tabla} & \textbf{Escalados} & \textbf{Espera fila} &
\textbf{Espera pág.} & \textbf{Total} \\
\midrule
entEMPRESA        &  4.066 & 83.106\,s & 55.258\,s & \critico{138.364\,s (38\,h)} \\
entORDENCARGA     & 16.398 & 54.744\,s & 25.453\,s & \critico{80.197\,s (22\,h)} \\
entALMACEN\_LOTE  &  1.840 &  5.367\,s & 20.023\,s & \critico{25.390\,s (7\,h)} \\
entPEDIDO         & 18.464 &  5.556\,s & 14.226\,s & 19.782\,s \\
entPEDIDO\_DETALLE& 14.044 &  6.938\,s &  7.033\,s & 13.971\,s \\
entMOVIMIENTO     & 11.281 &  6.232\,s &  2.845\,s &  9.077\,s \\
entOC\_DETALLE    &  2.761 &  2.211\,s &  6.778\,s &  8.989\,s \\
T\_Palets         & \aviso{193.154} & 551\,s & 4.259\,s & 4.810\,s \\
StnEDIPEDIDO      & 15.832 & 751\,s & 1.092\,s & 1.843\,s \\
STN\_720\_DATOS\_PESADA & 17.140 & 258\,s & 1.835\,s & 2.093\,s \\
APP\_PA\_StockArticulo & \aviso{68.447} & 143\,s & 0\,s & 143\,s \\
\bottomrule
\end{tabular}
\end{table}

\paragraph{El número de escalados engaña; el tiempo de espera no.}
\cod{T\_Palets} escala 193.154 veces y apenas genera espera.
\cod{entEMPRESA} escala solo 4.066 veces y acumula \textbf{38 horas de espera
por bloqueos}. Lo que importa no es cuántas veces se escala, sino sobre qué
tabla y con cuánta concurrencia alrededor.

\paragraph{Y \cod{entEMPRESA} tiene nombre y apellidos en la lista de jobs.}
Tres procesos recurrentes trabajan sobre esa tabla en pleno horario laboral:

\begin{itemize}\itemsep1pt
  \item \cod{BLOQUEAR EMPRESAS CON IMPAGOS} — cada hora, de 07:10 a 19:10
  \item \cod{BLOQUEAR EMPRESAS CON IMPAGOS MAIL} — cada hora, de 07:15 a 18:15
  \item \cod{DesbloquearEmpresasSaldoNegativo} — cada 30 min, de 08:00 a 19:00
\end{itemize}

Una tabla maestra de empresas, que casi cualquier consulta del ERP necesita
leer, reescrita cada media hora durante toda la jornada, sin versionado de
filas. \textbf{Es la fuente de bloqueo concreta más importante localizada.}

\paragraph{Hay una medida intermedia antes de RCSI.}
\cod{ALTER TABLE \ldots SET (LOCK\_ESCALATION = DISABLE)} sobre esas tres
tablas impide la conversión a bloqueo de tabla, a costa de más memoria de
bloqueos. Es quirúrgica y reversible en segundos, por lo que conviene
probarla antes que el cambio de aislamiento.

\subsection{Ninguna base de datos usa versionado de filas}

\begin{table}[h]
\centering
\small
\begin{tabular}{lcccc}
\toprule
\textbf{Base de datos} & \textbf{Compat.} & \textbf{RCSI} &
\textbf{Snapshot} & \textbf{Stats async} \\
\midrule
Grupo         & 160 & \critico{false} & \critico{OFF} & \critico{false} \\
Cenusa        & 160 & \critico{false} & \critico{OFF} & \critico{false} \\
OMS           & 160 & \critico{false} & \critico{OFF} & \critico{false} \\
IFR           & 160 & \critico{false} & \critico{OFF} & \critico{false} \\
ContaSTN      & 160 & \critico{false} & \critico{OFF} & \critico{false} \\
ContaAlaplana & 160 & \critico{false} & \critico{OFF} & \critico{false} \\
ContaVitaCer  & 160 & \critico{false} & \critico{OFF} & \critico{false} \\
D365Eventos   & 160 & \critico{false} & \critico{OFF} & \critico{false} \\
Informatica   & 160 & \critico{false} & \critico{OFF} & \critico{false} \\
\bottomrule
\end{tabular}
\end{table}

\cod{READ\_COMMITTED\_SNAPSHOT} está desactivado en todas, y el aislamiento
\emph{snapshot} también. Toda lectura toma bloqueos compartidos: \textbf{los
lectores bloquean a los escritores y viceversa}. Con 200 sesiones de usuario
leyendo mientras 140 jobs escriben, es exactamente la receta del cuadro
observado.

Activar RCSI en \cod{Grupo} haría que los lectores usaran versiones de fila
en lugar de bloqueos compartidos, y la mayor parte del bloqueo
lector-escritor desaparecería. \textbf{Es el cambio con mayor efecto
potencial sobre el síntoma que dispara la alarma.} Contrapartidas: consume
\emph{version store} en TempDB (ver \ref{sec:logs}, donde se demuestra que
hay margen de sobra) y cambia la semántica de lectura, por lo que exige
validación con la aplicación y una ventana con acceso exclusivo.

\paragraph{Lo positivo.} El nivel de compatibilidad es 160 en todas y
\cod{LEGACY\_CARDINALITY\_ESTIMATION} está en \cod{false}: el estimador
moderno ya está en uso, no hay nada que arreglar ahí.

\subsection{La carga real del SQL Agent}

El primer indicio fue una foto instantánea: quince jobs corriendo a la vez,
ocho de ellos arrancados \textbf{en el mismo segundo} (11:10:02), en su
mayoría sincronizaciones con Dynamics 365 y con SCALE. Al consultar
\cod{msdb.dbo.sysschedules} apareció el cuadro completo:
\critico{140 planificaciones sub-diarias activas}.

\begin{table}[h]
\centering
\small
\begin{tabular}{lrlrlr}
\toprule
\textbf{Cada} & \textbf{Jobs} & \textbf{Cada} & \textbf{Jobs} &
\textbf{Cada} & \textbf{Jobs} \\
\midrule
10\,s  & 3 & 3\,min  & 10 & 30\,min & 10 \\
30\,s  & 5 & 4\,min  & 4  & 1\,h    & 24 \\
75--100\,s & 5 & 5\,min & 13 & 2--4\,h & 15 \\
1\,min & 6 & 10\,min & 6  & 6--12\,h & 17 \\
2\,min & 7 & 15--18\,min & 10 & otros & 5 \\
\bottomrule
\end{tabular}
\caption{Seis jobs cada minuto, siete cada dos, trece cada cinco. Tres
\emph{cada diez segundos}. Solo los veinte más frecuentes suman unas 133.000
ejecuciones en siete días: 19.000 al día.}
\end{table}

\subsubsection*{Algunos tardan más que su propio intervalo}

\begin{table}[h]
\centering
\small
\begin{tabular}{lrrrr}
\toprule
\textbf{Job} & \textbf{Programado} & \textbf{Media} & \textbf{Máximo} &
\textbf{Ocupación} \\
\midrule
D365\_SINCRONIZAR\_OMS\_DISP\_MAX  & 120\,s & 142\,s & 671\,s  & \critico{118\,\%} \\
D365\_SYNC\_TRANSFORMAR\_DISP\_MAX & 180\,s & 173\,s & 689\,s  & \critico{96\,\%} \\
D365\_ENVIO\_RECEPCIONES           &  60\,s &  47\,s & 913\,s  & \critico{78\,\%} \\
stnSCALE\_COPIA\_SHIPPING\_LOAD    &  10\,s &   7\,s & 476\,s  & \critico{70\,\%} \\
D365\_SINCRONIZAR\_DATOS\_COMPRAS  & 300\,s & 200\,s & 707\,s  & \aviso{67\,\%} \\
D365\_SYNC\_MANTENER\_RESERVAS     & 240\,s &  98\,s & 667\,s  & 41\,\% \\
D365\_SYNC\_CREAR\_PEDIDO\_GRUPO   & 240\,s &  97\,s & 1.737\,s & 40\,\% \\
\bottomrule
\end{tabular}
\end{table}

\cod{D365\_SINCRONIZAR\_OMS\_DISP\_MAX} está programado cada dos minutos y
tarda de media dos minutos y veintidós segundos. \textbf{El «cada dos
minutos» es ficción}: SQL Agent no solapa una instancia consigo misma, así
que lo que ocurre es que arranca de nuevo en cuanto termina, de forma
continua, todo el día. \cod{D365\_SYNC\_TRANSFORMAR\_DISP\_MAX}, al 96\,\%,
está de hecho en la misma situación.

\subsubsection*{Los máximos delatan quién sufre el bloqueo}

\begin{table}[h]
\centering
\small
\begin{tabular}{lrrr}
\toprule
\textbf{Job} & \textbf{Media} & \textbf{Máximo} & \textbf{Factor} \\
\midrule
stnSCALE\_INT\_ENT\_PRODUCCION    &  5\,s & 1.781\,s (30\,min) & \critico{356$\times$} \\
Interface APP WEB                 & 16\,s & 4.590\,s (76\,min) & \critico{287$\times$} \\
STNAUTO\_AUTOMATIZAR\_ALBARANEADO &  8\,s & 1.161\,s (19\,min) & \critico{145$\times$} \\
STN\_MARCAR\_LINEAS\_PREPARADO    & 230\,s & 2.336\,s (39\,min) & 10$\times$ \\
\bottomrule
\end{tabular}
\end{table}

Un proceso que normalmente tarda cinco segundos y un día tarda treinta
minutos no se ha vuelto lento: se ha quedado esperando. Estos máximos
\textbf{no son causa, son consecuencia}, y demuestran que el bloqueo no solo
molesta a los usuarios: también atasca la propia automatización.

\subsubsection*{El lote nocturno invade la jornada}

\begin{table}[h]
\centering
\small
\begin{tabular}{llrl}
\toprule
\textbf{Job} & \textbf{Arranca} & \textbf{Máx.} & \textbf{Termina (peor caso)} \\
\midrule
RAPPEL\_ZCALCULOSFACTURAS & 04:00 & 11\,h 24\,m & \critico{15:24} \\
IndexOptimize USER\_DATABASES Weekend & 02:30 & 13\,h 59\,m & \critico{16:29} \\
TAREAS NOCTURNAS STN & 01:00 & 8\,h 14\,m & \critico{09:14} \\
STN\_MANTENER\_ARTICULOS\_CONSUMOS & --- & 3\,h 09\,m & --- \\
BAJA\_RIESGOS & 04:00 & 2\,h 16\,m & 06:16 \\
ACTUALIZAR FOTO CARTERA ARTICULO & 05:30 & 1\,h 24\,m & 06:54 \\
\bottomrule
\end{tabular}
\end{table}

\cod{TAREAS NOCTURNAS STN} puede terminar a las 09:14: \textbf{atraviesa de
lleno el pico de bloqueos de las 08:00}. El historial de \cod{IndexOptimize}
de los últimos ocho días confirma que el desbordamiento es real y recurrente:

\begin{table}[h]
\centering
\small
\begin{tabular}{llrl}
\toprule
\textbf{Fecha} & \textbf{Empieza} & \textbf{Duración} & \textbf{Termina} \\
\midrule
24/10 (sáb, \emph{Weekend}) & 02:30 & 13\,h 59\,m & \critico{16:29} \\
26/10 (dom, \emph{Daily})   & 02:30 &  6\,h 29\,m & \critico{08:59} \\
22/10 & 02:30 & 1\,h 25\,m & 03:55 \\
21, 23, 27, 28, 29/10 & 02:30 & $\approx$57\,min & $\approx$03:27 \\
25/10 & 02:30 & \multicolumn{2}{l}{\aviso{registro corrupto en el historial}} \\
\bottomrule
\end{tabular}
\end{table}

En días normales termina a las 03:27, pero el sábado llegó a las 16:29 y el
domingo a las 08:59. (La fila del 25/10 devuelve duración negativa y fecha de
fin en 2015: registro mal escrito en \cod{sysjobhistory}, probablemente de
una ejecución interrumpida.)

\subsubsection*{El equipo ya conocía esto}

Entre los jobs activos hay dos cuyos nombres son, en sí mismos, un hallazgo:

\begin{itemize}\itemsep1pt
  \item \cod{REVISAR BLOQUEOS DE 7:30 a 8:30} — con tres planificaciones, una
        de ellas cada 2 minutos entre las 06:00 y las 09:00.
  \item \cod{CONTROL JOBS SINFIN} — cada 15 minutos, todo el día.
\end{itemize}

Alguien del equipo ya había identificado tanto la ventana de bloqueos de las
08:00 como la tendencia de los jobs a quedarse colgados, y construyó
vigilancia para ambas cosas. \textbf{Este informe no descubre el problema: lo
cuantifica.}

\subsection{Hay 25\,GB de memoria sin usar}

\begin{table}[h]
\centering
\small
\begin{tabular}{lr}
\toprule
\textbf{Medida} & \textbf{Valor} \\
\midrule
RAM física del servidor        & 131.071\,MB (128\,GB) \\
Tope de SQL Server             & 96.896\,MB (94,6\,GB) \\
Lo que SQL Server está usando  & 96.917\,MB — su tope completo \\
RAM física libre               & \critico{25.657\,MB (25\,GB)} \\
\emph{Page life expectancy}    & \critico{28\,s} \\
\cod{system\_memory\_state\_desc} & \bien{Available physical memory is high} \\
\cod{process\_physical\_memory\_low} & \bien{false} \\
\bottomrule
\end{tabular}
\end{table}

Dejar 33\,GB al sistema operativo en un servidor de 128\,GB es conservador;
lo habitual sería reservar entre 16 y 20\,GB, lo que permitiría subir el tope
unos 12--15\,GB. Y el sistema operativo confirma que esa RAM está realmente
disponible: \emph{«Available physical memory is high»}, sin presión de
memoria externa sobre el proceso. El nivel de hipervisor lo corrobora
(\ref{sec:hipervisor}).

\subsection{La CPU: ocupada, y es SQL Server quien la ocupa}

\begin{table}[h]
\centering
\small
\begin{tabular}{lr}
\toprule
\textbf{Medida (10 minutos, 11:35--11:44)} & \textbf{Valor} \\
\midrule
CPU de SQL Server, media   & \aviso{72\,\%} \\
CPU de SQL Server, máximo  & \critico{83\,\%} \\
CPU de otros procesos      & \bien{3--6\,\%} \\
CPU libre, mínimo          & \critico{12\,\%} \\
Tareas en cola (\emph{runnable})  & 5 \\
Trabajos sin hilo asignado & \bien{0} \\
\emph{Load factor} medio / máximo & 13 / 19 \\
\bottomrule
\end{tabular}
\end{table}

La buena noticia: no hay ningún proceso ajeno comiéndose la máquina
—antivirus, agentes de respaldo, otra instancia— y \cod{work\_queue\_count} =
0 descarta el agotamiento del grupo de hilos. La mala: 72\,\% de media con
picos del 83\,\% y solo un 12\,\% de margen, con 5 tareas esperando turno,
es una máquina al límite. Respalda el \cod{SOS\_SCHEDULER\_YIELD} del
7,8\,\%.

\subsection{El disco queda descartado}

Latencia real por fichero (\cod{dm\_io\_virtual\_file\_stats}, acumulada
desde el arranque). Los ocho ficheros con peor latencia de lectura de toda la
instancia —entre los que superan 10.000 lecturas— son \textbf{todos ficheros
de log}, y ninguno pasa de 9\,ms:

\begin{table}[h]
\centering
\small
\begin{tabular}{llrrr}
\toprule
\textbf{Base de datos} & \textbf{Fichero} & \textbf{ms/lect.} &
\textbf{ms/escr.} & \textbf{Escrituras} \\
\midrule
ContaAlaplana & IN3ContaTpl\_Log & 9 & 0 & 1.063.679 \\
D365Eventos   & D365Eventos\_log & 9 & 1 & --- \\
ContaVitaCer  & log              & 8 & 1 & --- \\
Informatica   & Informatica\_Log & 8 & \aviso{12} & --- \\
OMS           & OMS\_log         & 8 & 1 & --- \\
ContaSTN      & IN3ContaTpl\_Log & 7 & 0 & 3.386.531 \\
Grupo         & In3Tplnew\_log   & 7 & 2 & 281.926.025 \\
IFR           & IFR\_log         & 7 & 2 & 363.247.671 \\
\bottomrule
\end{tabular}
\end{table}

Siete a nueve milisegundos por lectura de log es normal, y las escrituras de
log —las que de verdad frenan a las transacciones— están en 0--2\,ms.
\textbf{El almacenamiento no contribuye al atasco.} Dos detalles menores:
\cod{Informatica\_Log} escribe a 12\,ms, con poco volumen; y los logs de
\cod{Grupo} e \cod{IFR} acumulan 9,5 y 22 millones de \emph{lecturas}, lo que
es esperable al estar en un grupo de disponibilidad, porque la réplica lee el
log continuamente.

\subsection{Un índice que falta destaca sobre todos}

\cod{sys.dm\_db\_missing\_index\_details} en \cod{Grupo}. El primer candidato
supera al segundo \textbf{56 veces}:

\begin{table}[h]
\centering
\small
\begin{tabular}{llrr}
\toprule
\textbf{BD} & \textbf{Tabla} & \textbf{Impacto estimado} & \textbf{Búsquedas} \\
\midrule
Grupo & \critico{cmCargaConfirmadaLogs} \cod{[TipoAccion]} &
  \critico{45.515.122.075} & \critico{14.889.907} \\
Grupo & entMOVIMIENTO    &   805.654.619 & 1.165.115 \\
Grupo & stnERPGDExpDoc \cod{[IdDocumento]} & 437.142.615 & 260.722 \\
Grupo & ExpExpedicion    &   422.831.787 &   141.554 \\
Grupo & entALBARAN \cod{[atrIdLibro]} & 379.548.793 & 2.762 \\
Grupo & DELETED\_RECID   &   334.384.162 &    19.746 \\
Grupo & entALBARAN       &   287.714.517 &     9.724 \\
OMS   & PurchaseOrderHeader & (menor)    &       --- \\
\bottomrule
\end{tabular}
\end{table}

Casi quince millones de búsquedas en 27 días sobre una tabla sin ese índice
significa recorridos completos repetidos, y un recorrido completo es
exactamente lo que dispara un plan paralelo. Siete de los ocho candidatos
están en \cod{Grupo}.

\paragraph{Con una advertencia.} Un índice de una sola columna de baja
cardinalidad como \cod{[TipoAccion]} rara vez es el índice definitivo
correcto. La cifra sirve para saber \emph{dónde mirar}, no para aplicarla tal
cual: el paso siguiente es revisar las consultas reales que atacan esa tabla
y diseñar el índice a partir de ellas.

\subsection{El hipervisor queda descartado}
\label{sec:hipervisor}

Registrado el vCenter \cod{vcsa.central.local} en DPA, las cinco instancias
aparecen vinculadas a su máquina virtual. SRVPRORPDB01 corre sobre el host
ESXi \cod{10.150.0.8}, donde no comparte sitio con ninguna otra instancia
monitorizada (PROMETEO y SRVPBI sí comparten el \cod{10.150.0.9}).

\begin{table}[h]
\centering
\small
\begin{tabular}{llrr}
\toprule
\textbf{Métrica} & \textbf{Qué mide} & \textbf{Media} & \textbf{Máximo} \\
\midrule
\cod{vmReadyTime}        & espera por CPU física   & \bien{0,09\,\%} & \bien{0,13\,\%} \\
\cod{vmCpuCoStopPercent} & desincronía entre vCPU  & \bien{0\,\%}    & \bien{0\,\%} \\
\cod{vmPercentBallooned} & memoria reclamada       & \bien{0\,\%}    & \bien{0\,\%} \\
\cod{vmMemorySwapInRate} & intercambio del host    & \bien{0}        & \bien{0} \\
\cod{vmMemorySwapOutRate}& intercambio del host    & \bien{0}        & \bien{0} \\
\cod{vmMemoryGranted}    & memoria concedida       & 130.778\,MB     & 130.778\,MB \\
\cod{vmMemoryActive}     & conjunto activo         & 71,2\,\%        & 74,5\,\% \\
\cod{hostCpuUsagePercent}& CPU del host            & \bien{40,2\,\%} & \bien{45,5\,\%} \\
\cod{hostMemoryUsage}    & memoria del host        & \bien{39,4\,\%} & \bien{39,4\,\%} \\
\cod{vmCpuUsagePercent}  & CPU usada por la VM     & \aviso{96,4\,\%}& \aviso{114,8\,\%} \\
\bottomrule
\end{tabular}
\end{table}

\paragraph{Veredicto limpio.} El \emph{ready time} está en 0,09\,\%, cuando
el umbral habitual de preocupación ronda el 5\,\% y el de problema serio
supera el 10\,\%. El \emph{co-stop} —relevante precisamente en máquinas
anchas de 16 vCPU— es cero. \textbf{La máquina virtual no espera por CPU
física.} En memoria, las tres señales a descartar están a cero, con los
128\,GB concedidos íntegros y el host al 39\,\%. Esto \textbf{confirma que
los 25\,GB libres son reales} y respaldados por RAM física.

\paragraph{Dato útil para planificar.} La VM consume prácticamente toda su
asignación de CPU (96,4\,\%; los picos sobre 100\,\% se explican por el
cálculo sobre frecuencia nominal con \emph{turbo} activo). Con \emph{ready
time} a cero y el host al 40\,\%, no es que le falte CPU del host: es que
\textbf{está usando los 16 vCPU que tiene}. Ampliar vCPU es técnicamente
viable, pero sigue siendo la última acción del plan.

\paragraph{Una falsa alarma que conviene explicar.}
\cod{hostMaxTotalDiskLatency} marca 59\,ms de media y 211\,ms de pico. Al
desglosar por dispositivo, los siete discos de fibra del host están todos
\textbf{por debajo de 1,5\,ms}:

\begin{table}[h]
\centering
\small
\begin{tabular}{lrr}
\toprule
\textbf{Dispositivo (host 10.150.0.8)} & \textbf{ms lectura} & \textbf{ms escritura} \\
\midrule
naa.\ldots 1367656699695a40 & \bien{0,0} & \bien{0,0} \\
naa.\ldots 5f6d65661aaa6803 & \bien{0,0} & \bien{0,0} \\
naa.\ldots 2c6b6566ef2fd7fe & \bien{0,1} & \bien{0,9} \\
naa.\ldots 8f6f6566a55ec291 & \bien{0,5} & \bien{0,0} \\
naa.\ldots 8ec6916839933029 & \bien{0,6} & \bien{0,1} \\
naa.\ldots 086f65669cc587fa & \bien{0,8} & \bien{0,7} \\
naa.\ldots a1bf9168335813e0 & \bien{0,8} & \bien{1,3} \\
\bottomrule
\end{tabular}
\end{table}

\cod{maxTotalLatency} es el máximo entre \emph{todos} los dispositivos del
host, incluidos los que no son almacén de datos: un dispositivo local o una
ruta inactiva basta para dispararlo, y es un artefacto conocido en vSphere.
Los dispositivos reales están sanos, lo que \textbf{concuerda con la medición
desde dentro del invitado} y confirma por dos vías independientes que el
almacenamiento no es el problema.

\paragraph{Advertencia de ventana.} Solo hay $\approx$70 minutos de
histórico. Las magnitudes son inequívocas, pero el pico de las 08:00 no se ha
observado aún a este nivel.

\subsection{El grupo de disponibilidad queda descartado}

\begin{table}[h]
\centering
\small
\begin{tabular}{lllll}
\toprule
\textbf{Grupo} & \textbf{Réplica} & \textbf{Modo} & \textbf{Rol} &
\textbf{Salud} \\
\midrule
AG-ATENEA2 & SRVPRORPDB01 & \bien{ASYNCHRONOUS\_COMMIT} & PRIMARY & \bien{HEALTHY} \\
AG-ATENEA2 & SRVPRORPDB02 & \bien{ASYNCHRONOUS\_COMMIT} & SECONDARY & \bien{HEALTHY} \\
\bottomrule
\end{tabular}
\end{table}

Se sospechó que el AG estuviera en confirmación síncrona, lo que haría que
cada \emph{commit} esperase a la réplica secundaria y alargase la retención
de bloqueos. \textbf{No es el caso}: ambas réplicas están en
\cod{ASYNCHRONOUS\_COMMIT}, con retraso de 0 segundos y estado
\cod{HEALTHY}. Las colas de envío son de 60\,KB y las de rehacer, de 31\,KB
en \cod{Grupo}. El AG no contribuye al bloqueo.

Sí es coherente con el volumen de escritura: las tasas de rehacer son de
65.609 en \cod{Grupo} y 49.590 en \cod{IFR}, lo que explica los millones de
lecturas sobre los ficheros de log vistos antes.

\subsection{Registros de transacciones y \emph{version store}}
\label{sec:logs}

\begin{table}[h]
\centering
\small
\begin{tabular}{lrrrr}
\toprule
\textbf{Base de datos} & \textbf{VLFs} & \textbf{Activos} &
\textbf{Tamaño medio} & \textbf{Log total} \\
\midrule
Grupo  & \aviso{1.940} & 12 & 63,9\,MB & 121,0\,GB \\
IFR    & 127 & 6 & \aviso{1.465,7\,MB} & \critico{181,8\,GB} \\
Cenusa & 97 & \aviso{82 (85\,\%)} & 110,0\,MB & 10,4\,GB \\
\bottomrule
\end{tabular}
\end{table}

Ninguno es grave, pero los tres fallan por motivos opuestos. \cod{Grupo}
tiene \textbf{1.940 VLFs}, por encima del millar que se toma como
referencia: no afecta al trabajo diario, pero ralentiza arranque,
recuperación y copias de log. \cod{IFR} tiene el problema contrario —127 VLFs
de \textbf{1,4\,GB cada uno}—: con fragmentos tan grandes, un solo VLF activo
inmoviliza 1,4\,GB de log. Y \cod{Cenusa} tiene el \textbf{85\,\% de sus VLFs
activos}.

Al margen de los VLFs, llama la atención el \textbf{tamaño absoluto de los
logs}: 181,8\,GB en \cod{IFR} y 121\,GB en \cod{Grupo}. Son coherentes con
los 363 y 281 millones de escrituras medidas, pero un log de ese tamaño
sugiere que en algún momento no se truncó, y conviene verificar la política
de copias de log.

\paragraph{Línea base del \emph{version store}, antes de tocar RCSI.}

\begin{table}[h]
\centering
\small
\begin{tabular}{lr}
\toprule
\textbf{TempDB, uso actual} & \textbf{Valor} \\
\midrule
\emph{Version store}   & \bien{31\,MB} \\
Objetos de usuario     & 181\,MB \\
Objetos internos       & 1.268\,MB \\
Espacio libre          & \bien{430.726\,MB (420\,GB)} \\
\bottomrule
\end{tabular}
\end{table}

El \emph{version store} ocupa hoy \textbf{31\,MB}, lo lógico con RCSI
desactivado. Es la cifra de partida contra la que medir después del cambio. Y
el dato relevante: hay \textbf{420\,GB libres} en TempDB. Aunque activar RCSI
en \cod{Grupo} multiplicara el \emph{version store} por mil, seguiría
sobrando sitio. \textbf{Esto reduce considerablemente el riesgo de la acción
9}: el argumento en contra ya no puede ser el espacio de TempDB, solo la
validación funcional.

\subsection{El bloqueo de esquema de 392 segundos}
\label{sec:schm}

Se propuso que lo causara \cod{IndexOptimize}, ya que una reconstrucción de
índice toma un bloqueo \cod{Sch-M} incompatible con cualquier lectura.
\textbf{Los horarios lo desmienten}: el 29 de octubre ese job se ejecutó de
02:30 a 03:26:58, siete horas antes del incidente de la mañana. Hipótesis
descartada.

La explicación más probable ahora es otra: un \cod{TRUNCATE TABLE} sobre
alguna de las tablas de paso que se vacían y recargan continuamente
—\cod{entEFECTO\_FacturasCobradas365} es la candidata evidente—, porque
\cod{TRUNCATE} también toma \cod{Sch-M}. \textbf{Es una hipótesis, no un
hecho}: confirmarla requiere capturar el \emph{blocked process report} o una
sesión de \emph{Extended Events}, que no se ha hecho.

% =====================================================================
\newpage
\section{La base de datos \cod{Cenusa}}

Esta base de datos apareció por casualidad, al ordenar las consultas por
memoria concedida. No estaba en el foco inicial —el análisis se centraba en
\cod{Grupo}, que es donde se veían las cadenas de bloqueo— pero vive en la
misma instancia, ocupa 46\,GB y concentra \textbf{los costes unitarios más
altos de toda la investigación}.

\begin{table}[h]
\centering
\small
\begin{tabular}{lr}
\toprule
\textbf{Característica} & \textbf{Valor} \\
\midrule
Tamaño                              & 47.270\,MB (46\,GB) \\
Nivel de compatibilidad             & 160 \\
Estado                              & ONLINE \\
\cod{READ\_COMMITTED\_SNAPSHOT}     & \critico{false} \\
\cod{AUTO\_UPDATE\_STATISTICS\_ASYNC} & \critico{false} \\
\textbf{Disparadores sobre tablas}  & \critico{575} \\
De ellos, deshabilitados            & 4 \\
\bottomrule
\end{tabular}
\end{table}

\subsection{Quinientos setenta y cinco disparadores}

Solo cuatro están desactivados. Y están concentrados en las tablas más
transitadas:

\begin{table}[h]
\centering
\small
\begin{tabular}{lrr}
\toprule
\textbf{Tabla} & \textbf{Disparadores} & \textbf{Modificaciones (27 días)} \\
\midrule
entEMPRESA           & \critico{9} &  8.759.807 \\
entALBARAN           & \critico{8} & 26.778.967 \\
entMOVIMIENTO        & \critico{8} &  5.007.080 \\
entPEDIDO            & \critico{8} & 24.107.314 \\
entPEDIDO\_DETALLE   & \critico{8} & \critico{37.738.378} \\
entARTICULO          & 6 & 11.772.885 \\
entOC\_DETALLE       & 6 & --- \\
entARTICULO\_UBICACION & 5 & --- \\
\bottomrule
\end{tabular}
\end{table}

\paragraph{Por qué esto importa para los bloqueos.} Un disparador
\textbf{se ejecuta dentro de la transacción de quien lo dispara}. No es
trabajo que ocurra «después»: es tiempo que la sentencia original pasa
reteniendo sus bloqueos. Cada \cod{UPDATE} sobre \cod{entPEDIDO\_DETALLE}
arrastra ocho disparadores antes de poder confirmar. Con 37,7 millones de
modificaciones sobre esa tabla, la amplificación es enorme.

\subsection{Lo que cuestan}

\begin{table}[h]
\centering
\small
\begin{tabular}{lrrrr}
\toprule
\textbf{Objeto} & \textbf{Ejecuciones} & \textbf{CPU} &
\textbf{Mem. concedida} & \textbf{Aprov.} \\
\midrule
\cod{ACTUALIZA\_EMPRESAS\_DIRECCIONES\ldots} & 336 & \critico{4.828\,s} &
  \critico{1.171\,MB} & \critico{17,4\,\%} \\
\cod{tU\_entALBARAN} & 21.365 & 2.927\,s & --- & --- \\
\cod{tU\_entALBARAN} & \critico{1.700.015} & 1.851\,s & 2\,MB & 40,6\,\% \\
\cod{tU\_entMOVIMIENTO} ($\times$6) & 212.512 c/u & \critico{10.268\,s} & 1\,MB &
  0--51\,\% \\
\cod{update entFACTURA} ($\times$2) & 432.729 c/u & --- & 58\,MB & \critico{1,5\,\%} \\
\bottomrule
\end{tabular}
\end{table}

\paragraph{Un gigabyte de memoria para aprovechar el 17\,\%.} El
procedimiento \cod{ACTUALIZA\_EMPRESAS\_DIRECCIONES\_FICHAS\_POR\_PROVEEDOR}
recibe \textbf{1.171\,MB de concesión en cada una de sus 336 ejecuciones} y
usa el 17,4\,\%. Además lee \textbf{500.544 páginas por ejecución}, unos
3,8\,GB recorridos cada vez. Es, por amplio margen, la sentencia más cara
localizada en toda la instancia, y supera al caso de los 506\,MB de
\ref{sec:prueba}.

\paragraph{Un disparador con 31.000 lecturas por disparo.} La primera entrada
de \cod{tU\_entALBARAN} se ejecuta 21.365 veces leyendo \textbf{31.135
páginas cada vez}: unos 665 millones de lecturas lógicas acumuladas solo por
ese disparador. Y la segunda entrada se dispara \textbf{1,7 millones de
veces}.

\paragraph{Las seis caras del mismo disparador.}
\cod{tU\_entMOVIMIENTO} aparece seis veces porque son seis sentencias
distintas dentro del mismo cuerpo, cada una ejecutada 212.512 veces. Entre
las seis suman \textbf{10.268 segundos de CPU}, casi tres horas, y varias
aprovechan el 0\,\% de la memoria concedida.

\paragraph{Y 58\,MB aprovechados al 1,5\,\%.} Dos sentencias
\cod{update entFACTURA set atrBaseBruta\ldots}, con el patrón característico
de un cuerpo de disparador, se ejecutan 432.729 veces cada una reservando
58\,MB y usando el 1,5\,\%.

\subsection{\cod{AUDITORIA2}: la tabla que lo explica}

\begin{table}[h]
\centering
\small
\begin{tabular}{lrrr}
\toprule
\textbf{Tabla} & \textbf{Filas} & \textbf{Modificaciones} & \textbf{Muestra} \\
\midrule
AUDITORIA2       & \critico{29.708.129} & \critico{89.149.199} & \critico{1,0\,\%} \\
entALBARAN\_IVA  & 20.384.782 &  81.387 & \critico{1,1\,\%} \\
entPEDIDO        &    920.167 & 167.803 & \aviso{3,7--4,4\,\%} \\
\bottomrule
\end{tabular}
\end{table}

Una tabla de auditoría de \textbf{29,7 millones de filas} que acumula
\textbf{89 millones de modificaciones} —tres veces su propio tamaño— desde
que se actualizaron sus estadísticas esa madrugada. Y muestreada al
\textbf{1\,\%}.

Esto cierra el círculo de \cod{Cenusa}: los 575 disparadores escriben en
\cod{AUDITORIA2} en cada modificación. Cada \cod{UPDATE} de la aplicación
arrastra sus ocho o nueve disparadores, que a su vez escriben auditoría,
dentro de la misma transacción y reteniendo los mismos bloqueos. El coste no
es solo de CPU: es de \textbf{duración de transacción}, que es exactamente lo
que genera el bloqueo.

El muestreo al 1\,\% en tablas de 20 y 30 millones de filas merece revisión
por sí mismo. Es el valor que aplica \cod{IndexOptimize} por defecto en
tablas grandes, y sobre esos volúmenes produce histogramas de poca utilidad.

\subsection{Esperas de bloqueo en \cod{Cenusa}}

\begin{table}[h]
\centering
\small
\begin{tabular}{lrrrr}
\toprule
\textbf{Tabla} & \textbf{Escalados} & \textbf{Espera fila} &
\textbf{Espera pág.} & \textbf{Modificaciones} \\
\midrule
entPEDIDO\_DETALLE   &  4.209 &   159\,s & \critico{10.792\,s} & 37.738.378 \\
entALBARAN           &     21 &   398\,s & 3.113\,s & 26.778.967 \\
entARTICULO          &    117 & 2.884\,s &   146\,s & 11.772.885 \\
entEMPRESA           &      0 & 1.465\,s &   361\,s &  8.759.807 \\
entEMBALAJE\_ARTICULO&  1.789 & 1.328\,s &    41\,s & 11.528.414 \\
entPEDIDO            & \aviso{69.119} &  77\,s & 1.314\,s & 24.107.314 \\
entMOVIMIENTO        & \aviso{52.554} &   0\,s & 1.264\,s &  5.007.080 \\
entFICHA             &  5.961 &    86\,s &   421\,s &     81.756 \\
\bottomrule
\end{tabular}
\end{table}

Las cifras son un orden de magnitud menores que las de \cod{Grupo}
(\ref{sec:escalado}), donde \cod{entEMPRESA} acumula 138.364\,s. Pero el
volumen de modificaciones es enorme: 37,7 millones sobre
\cod{entPEDIDO\_DETALLE} y 26,8 sobre \cod{entALBARAN}, cada una con ocho
disparadores encima.

\subsection{Lo que sí está bien en \cod{Cenusa}}

\textbf{Los índices ausentes no son un problema aquí.} El mejor candidato
tiene un impacto estimado de 31 millones, frente a los 45.515 millones del
candidato de cabeza en \cod{Grupo}: unas \textbf{1.500 veces menos}.

\begin{table}[h]
\centering
\small
\begin{tabular}{llrr}
\toprule
\textbf{Tabla} & \textbf{Columnas} & \textbf{Impacto} & \textbf{Búsquedas} \\
\midrule
entMOVIMIENTO      & \cod{[atrIdTipoPedido]} & 30.973.289 & 8.298 \\
entPEDIDO          & \cod{[ifrCreadoDesdeD365]} & 30.497.282 & 9.142 \\
entPEDIDO\_DETALLE & \cod{[atrIdEmbalaje]} & 19.345.998 & 4.573 \\
entMOVIMIENTO      & \cod{[atrReferencia]} + tipo & 13.478.693 & \aviso{129.980} \\
entPEDIDO\_DETALLE & \cod{[atrPendiente]} & 13.265.771 & 1.363 \\
entFICHA           & \cod{[atrFechaPrecio]} & 3.814.413 & 914 \\
entARTICULO        & \cod{[atrArticuloActivo]} & 2.625.373 & 2.491 \\
\bottomrule
\end{tabular}
\end{table}

En \cod{Cenusa} el problema son los disparadores y el muestreo de
estadísticas, \textbf{no la falta de índices}.

\subsection{Qué hacer con esto}

Conviene decirlo sin adornos: \textbf{575 disparadores son una decisión de
arquitectura, no un parámetro}. Revisarlos uno a uno es trabajo de
desarrollo, de meses. Pero hay dos pasos acotados y de alto rendimiento:

\begin{enumerate}[leftmargin=1.4em]\itemsep2pt
  \item \textbf{Revisar
        \cod{ACTUALIZA\_EMPRESAS\_DIRECCIONES\_FICHAS\_POR\_PROVEEDOR}.} Una
        sola sentencia, 336 ejecuciones, 1,1\,GB de concesión y 3,8\,GB
        leídos por vez. Es el mejor retorno por hora invertida de todo el
        informe.
  \item \textbf{Revisar los disparadores de \cod{entMOVIMIENTO} y
        \cod{entALBARAN}}, que son los que más CPU consumen y los que más
        alargan las transacciones ajenas.
\end{enumerate}

\paragraph{Advertencia de alcance.} Este apartado abre un frente que
\textbf{no estaba en el encargo}. Las cifras son sólidas, pero no se ha
revisado el contenido funcional de los disparadores, y las decisiones sobre
cuáles tocar requieren conocimiento de la aplicación que este informe no
tiene. \textbf{Dado lo que ha aparecido, \cod{Cenusa} merece un análisis
propio.}

% =====================================================================
\newpage
\section{Hipótesis descartadas, y con qué se descartaron}
\label{sec:descartes}

Tan importante como lo que se recomienda es lo que \emph{no} hay que tocar.
Todo lo de esta tabla se comprobó por medición directa, no por criterio.

\begin{longtable}{p{4.1cm}p{5.1cm}p{5.4cm}}
\toprule
\textbf{Hipótesis} & \textbf{Qué se midió} & \textbf{Resultado} \\
\midrule
\endhead
Paralelismo mal configurado (instancia)
  & \cod{sys.configurations}
  & \bien{Descartado.} MAXDOP 8 sobre 16 núcleos y \emph{cost threshold} 100,
    frente al 5 por defecto \\
Paralelismo sobrescrito por base de datos
  & \cod{sys.database\_scoped\_configu\-rations} en \cod{Grupo}
  & \bien{Descartado.} MAXDOP\,=\,0 con \cod{is\_value\_default}\,=\,\cod{true};
    en ese ámbito el 0 significa «heredar la instancia» \\
Estimador de cardinalidad antiguo
  & \cod{compatibility\_level} y \cod{LEGACY\_CARDINALITY\_ESTIMATION}
  & \bien{Descartado.} Nivel 160 y \cod{false} en las nueve bases \\
TempDB mal dimensionado
  & \cod{tempdb.sys.database\_files}, \cod{dm\_db\_file\_space\_usage}
  & \bien{Descartado.} 8 ficheros de datos, 99\,\% libre, log al 93\,\% libre \\
Almacenamiento lento (invitado)
  & \cod{dm\_io\_virtual\_file\_stats} por fichero
  & \bien{Descartado.} Peor latencia de lectura 9\,ms; escrituras de log
    0--2\,ms \\
Almacenamiento lento (host)
  & \cod{hostTotalReadLatency}, \cod{hostTotalWriteLatency} por dispositivo
  & \bien{Descartado.} Los siete discos de fibra por debajo de 1,5\,ms \\
Proceso ajeno consumiendo CPU
  & \cod{RING\_BUFFER\_SCHEDULER\_MONITOR}
  & \bien{Descartado.} Los procesos ajenos a SQL Server, entre el 3 y el
    6\,\% \\
Agotamiento del grupo de hilos
  & \cod{dm\_os\_schedulers}
  & \bien{Descartado.} \cod{work\_queue\_count}\,=\,0 en todos \\
Contención de CPU en el hipervisor
  & \cod{vmReadyTime}, \cod{vmCpuCoStopPercent}
  & \bien{Descartado.} 0,09\,\% y 0\,\%. La VM no espera por CPU física \\
Presión de memoria del hipervisor
  & \cod{vmPercentBallooned}, \cod{vmMemorySwap*}
  & \bien{Descartado.} Todo a cero, 128\,GB concedidos \\
Confirmación síncrona del grupo de disponibilidad
  & \cod{sys.availability\_replicas}, \cod{dm\_hadr\_database\_replica\_states}
  & \bien{Descartado.} \cod{ASYNCHRONOUS\_COMMIT} en ambas réplicas, salud
    \cod{HEALTHY}, retraso 0\,s \\
Índices ausentes en \cod{Cenusa}
  & \cod{dm\_db\_missing\_index\_*}
  & \bien{Descartado.} Impacto 1.500 veces menor que en \cod{Grupo} \\
\cod{IndexOptimize} como causa del \cod{Sch-M} de 392\,s
  & \cod{msdb.dbo.sysjobhistory}
  & \bien{Descartado.} Terminó a las 03:26:58, siete horas antes \\
\bottomrule
\end{longtable}

\paragraph{Lo que no se recomienda hacer.}

\begin{itemize}[leftmargin=1.4em]\itemsep3pt
  \item \textbf{Matar sesiones de forma rutinaria.} Resuelve el síntoma del
        momento y oculta el problema. Solo tiene sentido como medida puntual
        ante una cadena parada varios minutos.
  \item \textbf{Tocar todo a la vez.} Si se cambian simultáneamente los jobs,
        la memoria y las estadísticas, no se sabrá cuál funcionó ni cuál
        empeoró algo.
  \item \textbf{Comprar vCPU antes de corregir las causas.} Con parte del
        72\,\% de CPU gastándose en paralelismo innecesario, es pagar por
        hacer lo mismo más rápido.
  \item \textbf{Aplicar los índices sugeridos tal cual.} Las cifras de
        impacto son heurísticas y no consideran los índices existentes ni el
        coste en escritura.
\end{itemize}

% =====================================================================
\newpage
\section{Plan de acción}
\label{sec:plan}

Catorce acciones en tres fases, cada una con su métrica de verificación.
Ordenadas por relación entre esfuerzo, riesgo y efecto.

\begin{longtable}{p{0.7cm}p{4.9cm}p{4.3cm}p{4.6cm}}
\toprule
\textbf{\#} & \textbf{Acción} & \textbf{Métrica de verificación} &
\textbf{Cómo saber que funcionó} \\
\midrule
\endhead
\multicolumn{4}{l}{\cellcolor{acento!10}\textbf{Fase 1 — esta semana.} Sin
parada, reversible, sin desarrollo.} \\
\midrule
1 & \cod{AUTO\_UPDATE\_STATISTICS\_ASYNC ON} en \cod{Grupo} y \cod{Cenusa}
  & esperas de compilación; \cod{SQL Compilations/sec}
  & desaparecen los picos de latencia sin causa aparente \\
2 & Estadísticas \cod{FULLSCAN} 3--4 veces al día sobre
    \cod{entEFECTO\_FacturasCobradas365}, \cod{STN\_ERP\_ANALISIS\_VENTAS\_PEDIDOS}
    y \cod{AUDITORIA2}
  & \cod{modification\_counter} de esas tablas
  & se mantiene por debajo del umbral, no en 26$\times$ el tamaño \\
3 & Alargar el intervalo de \cod{D365\_SINCRONIZAR\_OMS\_DISP\_MAX} y
    \cod{D365\_SYNC\_TRANSFORMAR\_DISP\_MAX}
  & \cod{pct\_ocupacion} (consulta \ref{sql:ocupacion})
  & por debajo del 50\,\% ambos \\
4 & \cod{LOCK\_ESCALATION = DISABLE} en \cod{entEMPRESA},
    \cod{entORDENCARGA} y \cod{entALMACEN\_LOTE}
  & \cod{row\_lock\_wait\_in\_ms} de esas tres tablas
  & caída clara frente a las 38, 22 y 7 horas actuales \\
5 & \cod{max server memory} de 94,6 a 107\,GB
  & \emph{Page life expectancy}
  & de 124\,s a más de 300\,s de media \\
\midrule
\multicolumn{4}{l}{\cellcolor{acento!10}\textbf{Fase 2 — este mes.} Requiere
análisis, sin cambiar la aplicación.} \\
\midrule
6 & Sacar el lote nocturno de la jornada: \cod{RAPPEL\_ZCALCULOSFACTURAS},
    \cod{TAREAS NOCTURNAS STN}, \cod{IndexOptimize}
  & hora de fin real de esos jobs
  & ninguno pasa de las 07:00 \\
7 & Escalonar los arranques coincidentes (anexo \ref{sec:anexo})
  & jobs activos por segundo
  & ningún segundo con más de 2 arranques \\
8 & Índice de \cod{Grupo.cmCargaConfirmadaLogs}, diseñado desde las consultas
    reales
  & impacto en \cod{dm\_db\_missing\_index\_*}
  & el candidato desaparece de la lista \\
9 & Parametrizar \cod{INSERT StnERPArticuloLoteLinea}
  & nº de planes distintos para esa sentencia
  & de 6 a 1; baja el CPU de 31.000\,s \\
10 & Activar Query Store en \cod{Grupo}
  & ---
  & habilita demostrar la evolución de los planes \\
11 & Revisar la política de copias de log y reducir los VLFs de \cod{Grupo}
  & \cod{dm\_db\_log\_info}
  & por debajo de 1.000 VLFs; logs que dejan de crecer \\
\midrule
\multicolumn{4}{l}{\cellcolor{acento!10}\textbf{Fase 3 — este trimestre.}
Requiere ventana, desarrollo o compra.} \\
\midrule
12 & \cod{READ\_COMMITTED\_SNAPSHOT ON} en \cod{Grupo}
  & sesiones bloqueadas; \cod{TempDB Version Store} (hoy 31\,MB, 420\,GB
    libres)
  & bloqueos casi a cero; el espacio de TempDB ya no es un riesgo \\
13 & \cod{Cenusa}: revisar
     \cod{ACTUALIZA\_EMPRESAS\_DIRECCIONES\ldots} y los disparadores de
     \cod{entMOVIMIENTO} y \cod{entALBARAN}
  & memoria concedida y lecturas lógicas por ejecución
  & bajar de 1.171\,MB y de 500.544 lecturas \\
14 & Valorar ampliar vCPU por encima de 16 —el host va al 40\,\%, hay holgura
  & \cod{SOS\_SCHEDULER\_YIELD}; \emph{runnable tasks}
  & solo si tras las fases 1--2 sigue alto \\
\bottomrule
\end{longtable}

\paragraph{Por qué este orden.} Las acciones 1 a 5 se pueden hacer esta misma
semana: no requieren parada y son reversibles. La 12 es la que más efecto
tendría sobre el síntoma que dispara la alarma, pero es la única que cambia
la semántica de lectura de la aplicación, así que va después de haber ganado
margen con las baratas —y ahora que se sabe que TempDB tiene 420\,GB libres,
su único riesgo real es funcional. La 14 se deja para el final a propósito.

\paragraph{Una regla para las tres fases.} \textbf{Un cambio a la vez, con 48
horas de observación entre uno y otro.} Si se tocan los jobs, la memoria y
las estadísticas el mismo día, el resultado será mejor o peor pero nadie
sabrá por qué, y no habrá forma de deshacer solo la parte que falló.

\subsection{Cómo medir si ha funcionado}

\begin{table}[h]
\centering
\small
\begin{tabular}{lrr}
\toprule
\textbf{Indicador} & \textbf{Hoy} & \textbf{Objetivo razonable} \\
\midrule
\% de muestras con bloqueo (7 días) & 88\,\% & por debajo del 30\,\% \\
Media de sesiones bloqueadas        & 4,9    & por debajo de 1 \\
Pico a las 8\,h                     & 9,0    & por debajo de 3 \\
\emph{Page life expectancy} medio   & 124\,s & por encima de 300\,s \\
Esperas \cod{CX*} sobre el total    & 58,7\,\% & por debajo del 30\,\% \\
Espera de bloqueo en \cod{entEMPRESA} & 138.364\,s & por debajo de 20.000\,s \\
\bottomrule
\end{tabular}
\end{table}

\subsection{Lo que queda pendiente de medir}

\begin{itemize}[leftmargin=1.4em]\itemsep3pt
  \item \textbf{Delta de esperas en ventana de 5 minutos durante el pico de
        las 08:00.} Todo el análisis de esperas es acumulado de 27 días; un
        delta en hora punta diría qué duele \emph{entonces}, no de media.
  \item \textbf{Métricas de hipervisor con 7 días completos}, para observar
        el pico de las 08:00 a ese nivel.
  \item \textbf{El origen real del bloqueo \cod{Sch-M}}, mediante
        \emph{blocked process report} o \emph{Extended Events}.
  \item \textbf{El contenido funcional de los 575 disparadores de
        \cod{Cenusa}.}
\end{itemize}

% =====================================================================
\newpage
\section{Consultas de diagnóstico}

Todas de solo lectura, ejecutadas sobre SRVPRORPDB01 con la cuenta
\cod{dpaMONITOR}. Reproducen punto por punto las cifras de este informe.

\subsection{Configuración de paralelismo, memoria y núcleos}

\begin{lstlisting}[style=sql]
SELECT name, value_in_use, description
FROM   sys.configurations
WHERE  name IN ('max degree of parallelism',
                'cost threshold for parallelism',
                'max server memory (MB)',
                'max worker threads');

SELECT cpu_count AS nucleos, hyperthread_ratio, softnuma_configuration_desc,
       sqlserver_start_time, virtual_machine_type_desc,
       scheduler_count AS schedulers,
       DATEDIFF(DAY, sqlserver_start_time, GETDATE()) AS dias_encendido
FROM   sys.dm_os_sys_info;
\end{lstlisting}

\subsection{Configuración por base de datos}

\begin{lstlisting}[style=sql]
SELECT name AS bd, compatibility_level AS compat,
       is_read_committed_snapshot_on   AS rcsi,
       snapshot_isolation_state_desc   AS snapshot_iso,
       is_auto_update_stats_on         AS auto_upd_stats,
       is_auto_update_stats_async_on   AS auto_upd_async,
       is_query_store_on               AS query_store
FROM   sys.databases
WHERE  database_id > 4
ORDER BY name;

USE Grupo;   -- repetir por cada base de interes

SELECT name AS opcion, value, value_for_secondary, is_value_default
FROM   sys.database_scoped_configurations
ORDER BY name;
\end{lstlisting}

\subsection{Esperas acumuladas, excluyendo las inactivas}

\begin{lstlisting}[style=sql]
SELECT TOP 12 wait_type,
       wait_time_ms/1000 AS espera_s,
       waiting_tasks_count AS veces,
       CAST(100.0*wait_time_ms/SUM(wait_time_ms) OVER() AS DECIMAL(5,2)) AS pct
FROM   sys.dm_os_wait_stats
WHERE  waiting_tasks_count > 0
  AND  wait_type NOT IN (
       'CLR_SEMAPHORE','LAZYWRITER_SLEEP','RESOURCE_QUEUE','SLEEP_TASK',
       'SLEEP_SYSTEMTASK','SQLTRACE_BUFFER_FLUSH','WAITFOR','LOGMGR_QUEUE',
       'CHECKPOINT_QUEUE','REQUEST_FOR_DEADLOCK_SEARCH','XE_TIMER_EVENT',
       'BROKER_TO_FLUSH','BROKER_TASK_STOP','CLR_MANUAL_EVENT',
       'DIRTY_PAGE_POLL','HADR_WORK_QUEUE','HADR_TIMER_TASK',
       'HADR_NOTIFICATION_DEQUEUE','HADR_FILESTREAM_IOMGR_IOCOMPLETION',
       'HADR_LOGCAPTURE_WAIT','HADR_CLUSAPI_CALL','SP_SERVER_DIAGNOSTICS_SLEEP',
       'QDS_ASYNC_QUEUE','QDS_SHUTDOWN_QUEUE','XE_DISPATCHER_WAIT')
ORDER BY wait_time_ms DESC;
\end{lstlisting}

\noindent La lista de exclusión es importante: sin ella,
\cod{HADR\_WORK\_QUEUE} por sí solo aporta 81 millones de segundos de espera
inactiva y distorsiona todos los porcentajes.

\subsection{Prueba del paralelismo innecesario}

\begin{lstlisting}[style=sql]
-- Consultas con plan paralelo que apenas devuelven filas
SELECT TOP 12 DB_NAME(st.dbid) AS bd, qs.execution_count AS ejec,
       qs.total_worker_time/1000000  AS cpu_total_s,
       qs.total_elapsed_time/1000000 AS dur_total_s,
       CAST(qs.total_worker_time*1.0/NULLIF(qs.total_elapsed_time,0)
            AS DECIMAL(6,2)) AS factor_par,
       qs.total_grant_kb/NULLIF(qs.execution_count,0)/1024 AS concedido_mb,
       CAST(100.0*qs.total_used_grant_kb/NULLIF(qs.total_grant_kb,0)
            AS DECIMAL(5,1)) AS pct_aprov,
       qs.total_rows/NULLIF(qs.execution_count,0) AS filas_media,
       LEFT(LTRIM(SUBSTRING(st.text,(qs.statement_start_offset/2)+1,
         ((CASE qs.statement_end_offset WHEN -1 THEN DATALENGTH(st.text)
           ELSE qs.statement_end_offset END - qs.statement_start_offset)/2)+1)),80)
         AS sentencia
FROM   sys.dm_exec_query_stats qs
CROSS APPLY sys.dm_exec_sql_text(qs.sql_handle) st
WHERE  qs.total_worker_time > qs.total_elapsed_time * 1.5   -- plan paralelo
  AND  qs.total_worker_time > 60000000                      -- > 60 s de CPU
ORDER BY qs.total_worker_time DESC;
\end{lstlisting}

\begin{lstlisting}[style=sql]
-- Memoria concedida frente a memoria usada: la huella de la mala estimacion
SELECT TOP 15 DB_NAME(st.dbid) AS bd, qs.execution_count AS ejec,
       qs.total_grant_kb/NULLIF(qs.execution_count,0)/1024      AS concedido_mb,
       qs.total_used_grant_kb/NULLIF(qs.execution_count,0)/1024 AS usado_mb,
       CAST(100.0*qs.total_used_grant_kb/NULLIF(qs.total_grant_kb,0)
            AS DECIMAL(5,1)) AS pct_aprovechado,
       qs.total_logical_reads/NULLIF(qs.execution_count,0) AS lecturas_media
FROM   sys.dm_exec_query_stats qs
CROSS APPLY sys.dm_exec_sql_text(qs.sql_handle) st
WHERE  qs.total_grant_kb > 0 AND qs.execution_count > 50
ORDER BY qs.total_grant_kb DESC;
\end{lstlisting}

\subsection{Obsolescencia de las estadísticas}

\begin{lstlisting}[style=sql]
USE Grupo;   -- o Cenusa

SELECT TOP 12 OBJECT_NAME(s.object_id) AS tabla, s.name AS estadistica,
       sp.last_updated, sp.rows AS filas,
       CAST(100.0*sp.rows_sampled/NULLIF(sp.rows,0) AS DECIMAL(5,1)) AS pct_muestra,
       sp.modification_counter AS modificaciones,
       DATEDIFF(DAY, sp.last_updated, GETDATE()) AS dias
FROM   sys.stats s
CROSS APPLY sys.dm_db_stats_properties(s.object_id, s.stats_id) sp
WHERE  sp.rows > 1000000
  AND  OBJECTPROPERTY(s.object_id, 'IsUserTable') = 1
ORDER BY sp.modification_counter DESC;
\end{lstlisting}

\subsection{Escalado y espera de bloqueos por tabla}

\begin{lstlisting}[style=sql]
USE Grupo;   -- o Cenusa

SELECT TOP 15 OBJECT_NAME(ios.object_id) AS tabla,
       SUM(ios.index_lock_promotion_count)         AS escalados,
       SUM(ios.index_lock_promotion_attempt_count) AS intentos,
       SUM(ios.row_lock_wait_count)                AS esperas_fila,
       SUM(ios.row_lock_wait_in_ms)/1000           AS espera_fila_s,
       SUM(ios.page_lock_wait_count)               AS esperas_pag,
       SUM(ios.page_lock_wait_in_ms)/1000          AS espera_pag_s,
       SUM(ios.leaf_insert_count + ios.leaf_update_count
           + ios.leaf_delete_count)                AS modificaciones
FROM   sys.dm_db_index_operational_stats(DB_ID(), NULL, NULL, NULL) ios
WHERE  OBJECTPROPERTY(ios.object_id,'IsUserTable') = 1
GROUP BY ios.object_id
HAVING SUM(ios.row_lock_wait_in_ms) + SUM(ios.page_lock_wait_in_ms) > 0
ORDER BY SUM(ios.row_lock_wait_in_ms) + SUM(ios.page_lock_wait_in_ms) DESC;
\end{lstlisting}

\subsection{Disparadores (base \cod{Cenusa})}

\begin{lstlisting}[style=sql]
USE Cenusa;

SELECT COUNT(*) AS triggers_tabla,
       SUM(CASE WHEN is_disabled = 1 THEN 1 ELSE 0 END) AS deshabilitados
FROM   sys.triggers WHERE parent_class = 1;

SELECT TOP 10 OBJECT_NAME(parent_id) AS tabla, COUNT(*) AS num_triggers
FROM   sys.triggers WHERE parent_class = 1 AND is_disabled = 0
GROUP BY parent_id ORDER BY COUNT(*) DESC;

-- Coste real, resolviendo el objeto que contiene cada sentencia
SELECT TOP 10 ISNULL(OBJECT_NAME(st.objectid, st.dbid),'(ad-hoc)') AS objeto,
       qs.execution_count AS ejec,
       qs.total_worker_time/1000000 AS cpu_s,
       qs.total_grant_kb/NULLIF(qs.execution_count,0)/1024 AS concedido_mb,
       CAST(100.0*qs.total_used_grant_kb/NULLIF(qs.total_grant_kb,0)
            AS DECIMAL(5,1)) AS pct_aprov,
       qs.total_logical_reads/NULLIF(qs.execution_count,0) AS lecturas_media
FROM   sys.dm_exec_query_stats qs
CROSS APPLY sys.dm_exec_sql_text(qs.sql_handle) st
WHERE  DB_NAME(st.dbid) = 'Cenusa'
ORDER BY qs.total_worker_time DESC;
\end{lstlisting}

\subsection{Índices ausentes}

\begin{lstlisting}[style=sql]
SELECT TOP 8
       DB_NAME(id.database_id) AS bd,
       OBJECT_NAME(id.object_id, id.database_id) AS tabla,
       CAST(igs.avg_total_user_cost * igs.avg_user_impact
            * (igs.user_seeks + igs.user_scans) AS BIGINT) AS impacto,
       igs.user_seeks, id.equality_columns, id.inequality_columns,
       id.included_columns
FROM   sys.dm_db_missing_index_group_stats igs
JOIN   sys.dm_db_missing_index_groups ig  ON ig.index_group_handle = igs.group_handle
JOIN   sys.dm_db_missing_index_details id ON id.index_handle = ig.index_handle
ORDER BY impacto DESC;
\end{lstlisting}

\subsection{Carga real del SQL Agent}
\label{sql:ocupacion}

\begin{lstlisting}[style=sql]
-- Cuantas planificaciones sub-diarias hay, y a que ritmo
SELECT CASE s.freq_subday_type WHEN 2 THEN 'segundos' WHEN 4 THEN 'minutos'
            WHEN 8 THEN 'horas' ELSE 'otro' END AS unidad,
       s.freq_subday_interval AS cada_n, COUNT(*) AS num_jobs
FROM   msdb.dbo.sysjobs j
JOIN   msdb.dbo.sysjobschedules js ON js.job_id = j.job_id
JOIN   msdb.dbo.sysschedules    s  ON s.schedule_id = js.schedule_id
WHERE  j.enabled = 1 AND s.enabled = 1 AND s.freq_subday_interval > 0
GROUP BY s.freq_subday_type, s.freq_subday_interval
ORDER BY unidad, cada_n;
\end{lstlisting}

\begin{lstlisting}[style=sql]
-- Jobs que tardan mas que su propio intervalo de planificacion
WITH prog AS (
  SELECT j.job_id, j.name,
         s.freq_subday_interval * CASE s.freq_subday_type
             WHEN 2 THEN 1 WHEN 4 THEN 60 WHEN 8 THEN 3600 ELSE 60 END AS intervalo_s
  FROM   msdb.dbo.sysjobs j
  JOIN   msdb.dbo.sysjobschedules js ON js.job_id = j.job_id
  JOIN   msdb.dbo.sysschedules    s  ON s.schedule_id = js.schedule_id
  WHERE  j.enabled = 1 AND s.enabled = 1 AND s.freq_subday_interval > 0
),
dur AS (
  SELECT h.job_id,
         AVG(h.run_duration/10000*3600 + (h.run_duration/100%100)*60
             + h.run_duration%100) AS media_s,
         MAX(h.run_duration/10000*3600 + (h.run_duration/100%100)*60
             + h.run_duration%100) AS max_s,
         COUNT(*) AS ejec
  FROM   msdb.dbo.sysjobhistory h
  WHERE  h.step_id = 0
    AND  h.run_date >= CONVERT(INT, CONVERT(VARCHAR(8), DATEADD(DAY,-7,GETDATE()), 112))
  GROUP BY h.job_id
)
SELECT TOP 20 p.name AS job, p.intervalo_s, d.media_s, d.max_s, d.ejec,
       CAST(100.0*d.media_s/NULLIF(p.intervalo_s,0) AS DECIMAL(8,1)) AS pct_ocupacion
FROM   prog p JOIN dur d ON d.job_id = p.job_id
WHERE  d.max_s > p.intervalo_s
ORDER BY pct_ocupacion DESC;
\end{lstlisting}

\begin{lstlisting}[style=sql]
-- Ejecuciones largas: cuando empezaron y cuando acabaron de verdad
SELECT TOP 25 j.name AS job,
       msdb.dbo.agent_datetime(h.run_date, h.run_time) AS empezo,
       h.run_duration/10000*3600 + (h.run_duration/100%100)*60
           + h.run_duration%100 AS dur_s,
       DATEADD(SECOND, h.run_duration/10000*3600 + (h.run_duration/100%100)*60
           + h.run_duration%100,
           msdb.dbo.agent_datetime(h.run_date, h.run_time)) AS acabo
FROM   msdb.dbo.sysjobhistory h
JOIN   msdb.dbo.sysjobs j ON j.job_id = h.job_id
WHERE  h.step_id = 0
  AND  h.run_date >= CONVERT(INT, CONVERT(VARCHAR(8), DATEADD(DAY,-8,GETDATE()), 112))
  AND  h.run_duration > 3000
ORDER BY dur_s DESC;
\end{lstlisting}

\subsection{Memoria, CPU y presión de hilos}

\begin{lstlisting}[style=sql]
SELECT total_physical_memory_kb/1024     AS ram_mb,
       available_physical_memory_kb/1024 AS libre_mb,
       system_memory_state_desc
FROM   sys.dm_os_sys_memory;

SELECT physical_memory_in_use_kb/1024 AS en_uso_mb,
       process_physical_memory_low, process_virtual_memory_low
FROM   sys.dm_os_process_memory;

SELECT SUM(runnable_tasks_count) AS tareas_en_cola,
       SUM(work_queue_count)     AS sin_hilo,
       SUM(current_tasks_count)  AS tareas_totales,
       SUM(active_workers_count) AS workers,
       MAX(load_factor) AS load_factor_max, AVG(load_factor) AS load_factor_medio
FROM   sys.dm_os_schedulers WHERE status = 'VISIBLE ONLINE';
\end{lstlisting}

\begin{lstlisting}[style=sql]
-- CPU de SQL Server frente a la del resto del sistema (ultimos ~256 min)
SELECT TOP 10
       DATEADD(ms, -1*((SELECT ms_ticks FROM sys.dm_os_sys_info) - y.[timestamp]),
               GETDATE()) AS hora,
       y.cpu_sql, 100 - y.cpu_libre - y.cpu_sql AS cpu_otros, y.cpu_libre
FROM (
  SELECT x.[timestamp],
         x.record.value('(./Record/SchedulerMonitorEvent/SystemHealth/SystemIdle)[1]',
                        'int') AS cpu_libre,
         x.record.value('(./Record/SchedulerMonitorEvent/SystemHealth/ProcessUtilization)[1]',
                        'int') AS cpu_sql
  FROM (SELECT [timestamp], CONVERT(XML, record) AS record
        FROM   sys.dm_os_ring_buffers
        WHERE  ring_buffer_type = 'RING_BUFFER_SCHEDULER_MONITOR'
          AND  record LIKE '%<SystemHealth>%') AS x
) AS y
ORDER BY y.[timestamp] DESC;
\end{lstlisting}

\subsection{Disco, logs y TempDB}

\begin{lstlisting}[style=sql]
SELECT TOP 8 DB_NAME(vfs.database_id) AS bd, mf.name AS fichero, mf.type_desc,
       vfs.num_of_reads,
       vfs.io_stall_read_ms  / NULLIF(vfs.num_of_reads , 0) AS ms_por_lectura,
       vfs.num_of_writes,
       vfs.io_stall_write_ms / NULLIF(vfs.num_of_writes, 0) AS ms_por_escritura
FROM   sys.dm_io_virtual_file_stats(NULL, NULL) AS vfs
JOIN   sys.master_files AS mf
       ON mf.database_id = vfs.database_id AND mf.file_id = vfs.file_id
WHERE  vfs.num_of_reads > 10000
ORDER BY ms_por_lectura DESC;

SELECT 'Grupo' AS bd, COUNT(*) AS vlfs,
       SUM(CASE WHEN vlf_active = 1 THEN 1 ELSE 0 END) AS activos,
       CAST(AVG(vlf_size_mb) AS DECIMAL(10,1)) AS tam_medio_mb,
       CAST(SUM(vlf_size_mb)/1024 AS DECIMAL(10,1)) AS log_total_gb
FROM   sys.dm_db_log_info(DB_ID('Grupo'));

SELECT SUM(version_store_reserved_page_count)*8/1024   AS version_store_mb,
       SUM(user_object_reserved_page_count)*8/1024     AS objetos_usuario_mb,
       SUM(internal_object_reserved_page_count)*8/1024 AS objetos_internos_mb,
       SUM(unallocated_extent_page_count)*8/1024       AS libre_mb
FROM   tempdb.sys.dm_db_file_space_usage;
\end{lstlisting}

\subsection{Grupo de disponibilidad}

\begin{lstlisting}[style=sql]
SELECT ag.name AS grupo, ar.replica_server_name AS replica,
       ar.availability_mode_desc AS modo_commit, ar.failover_mode_desc,
       ars.role_desc AS rol, ars.synchronization_health_desc AS salud
FROM   sys.availability_groups ag
JOIN   sys.availability_replicas ar ON ar.group_id = ag.group_id
LEFT JOIN sys.dm_hadr_availability_replica_states ars
       ON ars.replica_id = ar.replica_id;

SELECT DB_NAME(drs.database_id) AS bd, ar.replica_server_name AS replica,
       drs.synchronization_state_desc AS estado,
       drs.log_send_queue_size AS cola_envio_kb,
       drs.redo_queue_size AS cola_redo_kb, drs.redo_rate,
       drs.secondary_lag_seconds AS retraso_s
FROM   sys.dm_hadr_database_replica_states drs
JOIN   sys.availability_replicas ar ON ar.replica_id = drs.replica_id;
\end{lstlisting}

\subsection{Métricas de hipervisor (API interna de DPA)}

No son consultas SQL: DPA las obtiene del vCenter y las expone en su API
interna, la misma que consume la pestaña \emph{Virtualization}. Requieren
cookie de sesión web, \textbf{no} el token de la API pública documentada, que
es puramente administrativa y no expone métricas de rendimiento.

\begin{lstlisting}[style=sql]
# Instancias y su vinculacion con la maquina virtual
GET /iwc/rest/databases
    -> id, name, linkedToVm, virtualMachineName,
       virtualMachineCurrentHostName, vsphereServerName

# Catalogo de metricas en modo Virtualization (pm=V)
GET /iwc/rest/databases/{id}/metrics/categories
      ?startTime=...&endTime=...&pm=V&require=metrics,status

# Serie temporal de una metrica concreta
GET /iwc/rest/databases/{id}/metrics/data
      ?metrics=vmReadyTime&startTime=...&endTime=...
      &useResourceGranularities=true&baseline=false
    -> data[0].series[n].dataPoints[].{time,value}

# Metricas relevantes para descartar contencion del host:
#   CPU     : vmReadyTime, vmCpuCoStopPercent, hostCpuUsagePercent,
#             vmCpuUsagePercent, vmCpuUsed, vmCpuUsageMhz
#   Memoria : vmPercentBallooned, vmBalloonSize, vmMemorySwapInRate,
#             vmMemorySwapOutRate, vmMemoryGranted, vmMemoryActive,
#             hostMemoryUsage
#   Disco   : hostMaxTotalDiskLatency, hostTotalReadLatency,
#             hostTotalWriteLatency (estas dos, por dispositivo)
\end{lstlisting}

% =====================================================================
\newpage
\section{Evidencia}

Salida literal de la herramienta de consultas de DPA contra SRVPRORPDB01,
el 29 de octubre de 2026, conectada como \cod{dpaMONITOR}. Sin editar.

\subsection{Configuración}

\begin{lstlisting}[style=salida]
SELECT name, value_in_use FROM sys.configurations WHERE name IN (...)
Time: 94 ms

name                             value_in_use
cost threshold for parallelism   100
max degree of parallelism        8
max server memory (MB)           96896
\end{lstlisting}

\begin{lstlisting}[style=salida]
SELECT cpu_count, numa, page life expectancy, memoria ...
Time: 47 ms

nucleos   numa   ple_s   sql_usa_mb   ram_fisica_mb   ram_libre_mb
16        2      28      96896        131071          25657
\end{lstlisting}

\begin{lstlisting}[style=salida]
SELECT sqlserver_start_time, dias_encendido, virtual_machine_type_desc ...
Time: 16 ms

sqlserver_start_time      dias_encendido   virtual_machine_type_desc   schedulers
2026-10-02 14:25:05.69    27               HYPERVISOR                  16
\end{lstlisting}

\begin{lstlisting}[style=salida]
SELECT name AS bd, compatibility_level, is_read_committed_snapshot_on ...
Time: 140 ms

bd              compat  rcsi   snapshot_iso  auto_upd_stats  query_store  auto_upd_async
ContaAlaplana   160     false  OFF           true            false        false
ContaSTN        160     false  OFF           true            false        false
ContaVitaCer    160     false  OFF           true            false        false
D365Eventos     160     false  OFF           true            true         false
Grupo           160     false  OFF           true            false        false
IFR             160     false  OFF           true            false        false
Informatica     160     false  OFF           true            false        false
OMS             160     false  OFF           true            false        false

USE Grupo; SELECT name, value, is_value_default FROM sys.database_scoped_configurations
Time: 32 ms

opcion                                     value  value_for_secondary  is_value_default
ACCELERATED_PLAN_FORCING                   true   null                 true
ASYNC_STATS_UPDATE_WAIT_AT_LOW_PRIORITY    false  null                 true
LEGACY_CARDINALITY_ESTIMATION              false  null                 true
MAXDOP                                     0      null                 true
OPTIMIZE_FOR_AD_HOC_WORKLOADS              false  null                 true
PARAMETER_SNIFFING                         true   null                 true
QUERY_OPTIMIZER_HOTFIXES                   false  null                 true
\end{lstlisting}

\subsection{Esperas}

\begin{lstlisting}[style=salida]
SELECT TOP 6 wait_type, wait_time_ms/1000, pct FROM sys.dm_os_wait_stats ...
Time: 15 ms

wait_type                   espera_s     pct
CXCONSUMER                  45504691     27,55
CXSYNC_PORT                 31022696     18,79
CXPACKET                    20425851     12,37
SOS_SCHEDULER_YIELD         12906011      7,82
LATCH_EX                     9896832      5,99
HADR_NOTIFICATION_DEQUEUE    4645761      2,81   <- espera inactiva
\end{lstlisting}

\subsection{Paralelismo innecesario}

\begin{lstlisting}[style=salida]
SELECT TOP 12 ... WHERE total_worker_time > total_elapsed_time * 1.5
Time: 1670 ms

bd      ejec  cpu_s  dur_s  factor_par  conc_mb  pct_aprov  filas  sentencia
Grupo    358   5535   2167      2.55        71      55.4       0   insert into StnERPArticuloLoteLinea (...)
Grupo    358   5370   2043      2.63        71      55.4       0   insert into StnERPArticuloLoteLinea (...)
Grupo    358   5349   2017      2.65        71      55.4       0   insert into StnERPArticuloLoteLinea (...)
Grupo    357   4963   2094      2.37        70      55.7       0   insert into StnERPArticuloLoteLinea (...)
Grupo    358   4960   1986      2.50        70      55.7       0   insert into StnERPArticuloLoteLinea (...)
Grupo    358   4819   1859      2.59        70      55.7       0   insert into StnERPArticuloLoteLinea (...)
Grupo   1122   1340    376      3.57        33      49.2       0   update stnERPORDENPREPARACION set Estado = case ...
Grupo   1122    919    328      2.80        33      46.6       0   update stnERPORDENPREPARACION set Estado = 2 ...
null     591   1993    898      2.22        27      55.9       3   delete b from [D365Eventos]..STNAppSuiteInventTrans
null     597   1255    597      2.10        25      58.0       3   delete b from [D365Eventos]..STNAppSuiteInventTrans
null     595   1181    627      1.88        18      43.2       2   delete b from [D365Eventos]..InventTransOriginBi...
null     164    581    291      2.00       506      67.0       1   Select a.SourceKey FROM OMS..InventTransOrigin...
\end{lstlisting}

\begin{lstlisting}[style=salida]
SELECT TOP 15 ... ORDER BY total_grant_kb DESC
Time: ---

bd       ejec      conc_mb  usado_mb  pct_aprov  factor  filas  sentencia
Cenusa    432729       58       0        1.5      0.48     0    update entFACTURA set atrBaseBruta = ...
Cenusa    432728       58       0        1.5      0.49     0    update entFACTURA set atrBaseBruta = ...
Cenusa   1699520        2       0       40.6      0.48     0    update entFACTURA set atrCostefactura = ...
Cenusa   1722472        1       0        0.0      0.92     0    select @TestPendiente = entPEDIDO_DETALLE...
Grupo    5631854        1       0        0.0      0.45     0    UPDATE T SET modified_date_time = GETDATE()...
Grupo     127838       42      28       67.3      0.70 103528   INSERT INTO dbo.SHIPPING_LOAD_COPIA (...)
Grupo    5260024        1       0        0.0      0.45     0    update entALMACEN_LOTE set atrAlPendiente...
\end{lstlisting}

\subsection{Estadísticas}

\begin{lstlisting}[style=salida]
USE Grupo; SELECT TOP 12 ... FROM sys.stats CROSS APPLY dm_db_stats_properties
Time: 22094 ms

tabla                              last_updated          filas    pct_mu  modificaciones
entEFECTO_FacturasCobradas365      2026-10-29 02:41:30  1184877    11,8      31476588
entEFECTO_FacturasCobradas365      2026-10-29 02:41:28  1184877    11,8      21329427
STN_ERP_ANALISIS_VENTAS_PEDIDOS    2026-10-29 02:48:51  6460673     1,4      12940889
STN_ERP_ANALISIS_VENTAS_PEDIDOS    2026-10-29 02:48:58  6460673     1,4      12940889

USE Cenusa; (misma consulta)
Time: 10743 ms

AUDITORIA2                         2026-10-29 02:54:54 29708129     1,0      89149199
entALBARAN_IVA                     2026-10-29 02:55:15 20384782     1,1         81387
entPEDIDO                          2026-10-29 02:56:35   920167     4,4        167803
\end{lstlisting}

\subsection{Escalado de bloqueos}

\begin{lstlisting}[style=salida]
USE Grupo; SELECT TOP 15 ... FROM sys.dm_db_index_operational_stats
Time: 13047 ms

tabla                  escalados  intentos  esp_fila  esp_fila_s  esp_pag  esp_pag_s
entEMPRESA                  4066   1257751     16919       83106     7457      55258
entORDENCARGA              16398   1246542      8118       54744     3472      25453
entALMACEN_LOTE             1840    308481       788        5367     3474      20023
entPEDIDO                  18464   2110454      1356        5556     5314      14226
entPEDIDO_DETALLE          14044   2400993      2004        6938     2563       7033
entMOVIMIENTO              11281   1508432       170        6232      371       2845
entOC_DETALLE               2761   1535207       625        2211     1478       6778
T_Palets                  193154   1513608      7345         551     2012       4259
STN_720_DATOS_PESADA       17140     69221       421         258      452       1835
StnEDIPEDIDO               15832    276422      2010         751      813       1092
APP_PA_StockArticulo       68447   6753289       458         143        0          0
ENVIO_MAILS                42475   1520206       432         268       13          3
entFICHA                    5961      6842        56          86       45        421
\end{lstlisting}

\subsection{Cenusa}

\begin{lstlisting}[style=salida]
SELECT d.name, compatibility_level, rcsi, async_stats, state_desc, tamano_mb
Time: 94 ms

bd       compat  rcsi   async_stats  estado   tamano_mb
Cenusa   160     false  false        ONLINE   47270

SELECT COUNT(*) FROM sys.triggers WHERE parent_class = 1
Time: 15 ms

triggers_tabla   deshabilitados
575              4

SELECT TOP 8 OBJECT_NAME(parent_id), COUNT(*) ... GROUP BY parent_id
Time: 32 ms

tabla                    num_triggers
entEMPRESA                      9
entALBARAN                      8
entMOVIMIENTO                   8
entPEDIDO                       8
entPEDIDO_DETALLE               8
entARTICULO                     6
entOC_DETALLE                   6
entARTICULO_UBICACION           5
\end{lstlisting}

\begin{lstlisting}[style=salida]
SELECT TOP 10 ... WHERE DB_NAME(st.dbid) = 'Cenusa' ORDER BY total_worker_time DESC
Time: 344 ms

objeto                                           ejec     cpu_s  dur_s  conc_mb  pct   lecturas
ACTUALIZA_EMPRESAS_DIRECCIONES_FICHAS_POR_PROV     336      4828   4208     1171  17.4   500544
tU_entALBARAN                                    21365      2927   4829        0  null    31135
tU_entALBARAN                                  1700015      1851   3861        2  40.6        0
tU_entMOVIMIENTO                                212512      1884   4008        1  51.5       16
tU_entMOVIMIENTO                                212512      1860   3963        1  51.5       16
tU_entMOVIMIENTO                                212512      1757   3699        1   0.0       16
tU_entMOVIMIENTO                                212512      1691   3599        1   0.0       16
tU_entMOVIMIENTO                                212512      1588   4218        1   0.0       35
tU_entMOVIMIENTO                                212512      1488   3353        1   0.0       --
\end{lstlisting}

\subsection{SQL Agent}

\begin{lstlisting}[style=salida]
SELECT ... COUNT(*) FROM msdb.dbo.sysschedules WHERE freq_subday_interval > 0
Time: 141 ms

unidad      cada_n  jobs     unidad    cada_n  jobs     unidad   cada_n  jobs
segundos        10     3     minutos       10     6     horas         1    24
segundos        30     5     minutos       15     9     horas         2     8
segundos        75     1     minutos       18     1     horas         3     3
segundos        90     2     minutos       30    10     horas         4     4
segundos       100     2     minutos       60     1     horas      6-12    17
minutos          1     6                                           otro  1-10     4
minutos          2     7                                           ------------------
minutos          3    10                                           TOTAL          140
minutos          4     4
minutos          5    13
\end{lstlisting}

\begin{lstlisting}[style=salida]
WITH prog AS (...), dur AS (...) SELECT TOP 20 ... WHERE d.max_s > p.intervalo_s
Time: 2484 ms

job                                     intervalo_s  media_s  max_s  ejec    pct
D365_SINCRONIZAR_OMS_DISP_MAX                   120      142    671   2380  118,3
D365_SYNC_TRANSFORMAR_DISP_MAX                  180      173    689   1977   96,1
D365_ENVIO_RECEPCIONES                           60       47    913   6202   78,3
stnSCALE_COPIA_SHIPPING_LOAD                     10        7    476  34670   70,0
D365_SINCRONIZAR_DATOS_COMPRAS                  300      200    707   1215   66,7
D365_SYNC_MANTENER_RESERVAS_GRUPO_DESDED365     240       98    667   2532   40,8
D365_SYNC_CREAR_PEDIDO_GRUPO_DESDED365          240       97   1737   2600   40,4
STN_MARCAR_LINEAS_PREPARADO_TRASLADADO          600      230   2336   1043   38,3
STN_SINCRONIZASTOCKRUMBO                        300      114    718   1433   38,0
stnSCALE_720_DATOS_PESADA_ACTUALIZA_MUELLE       10        3     80  29830   30,0
D365_SYNC_CREAR_OC_GRUPO_DESDE_D365             240       68    641   2676   28,3
stnSCALE_INT_ENT_PRODUCCION                      30        5   1781  19432   16,7
Interface APP WEB                               120       16   4590   5346   13,3
STNAUTO_AUTOMATIZAR_ALBARANEADO                  60        8   1161  10487   13,3
\end{lstlisting}

\begin{lstlisting}[style=salida]
SELECT TOP 20 j.name, AVG(run_duration), MAX(run_duration) ... (7 dias)
Time: 1158 ms

job                                        ejec  dur_media_s  dur_max_s  fallos
IndexOptimize - USER_DATABASES - Weekend      1        50360      50360       0
RAPPEL_ZCALCULOSFACTURAS                      7         8890      41020       0
TAREAS NOCTURNAS STN                          8         8801      29672       0
STN_MANTENER_ARTICULOS_CONSUMOS               7         8217      11323       0
RECOSTRUYE STOCK DOMINGO                      1         7079       7079       0
BAJA_RIESGOS                                  8         6766       8129       0
STN_ERP_CARGAR_ANALISIS_VENTAS_HISTORICO      8         3121       3960       0
STN_AVISO_DISPONIBLE_FASE                    32         2609       6065       0
IN3 RUMBO INVENTARIO                          7         2259       3136       0
DatabaseBackup - GRUPO_DATABASES - FULL      15         2202       2592       0
D365_SYNC_BORRAR_OMS                         96         1133       1466       2
ACTUALIZAR FOTO CARTERA ARTICULO              6         1072       5067       0
\end{lstlisting}

\begin{lstlisting}[style=salida]
SELECT ... WHERE j.name LIKE '%IndexOptimize%' ... (8 dias)
Time: 4810 ms

job                                        empezo                 dur_s   acabo
IndexOptimize - USER_DATABASES - Daily     2026-10-29 02:30:00     3418   2026-10-29 03:26:58
IndexOptimize - USER_DATABASES - Daily     2026-10-28 02:30:01     3403   2026-10-28 03:26:44
IndexOptimize - USER_DATABASES - Daily     2026-10-27 02:30:00     3498   2026-10-27 03:28:18
IndexOptimize - USER_DATABASES - Daily     2026-10-26 02:30:01    23366   2026-10-26 08:59:27  <--
IndexOptimize - USER_DATABASES - Daily     2026-10-25 02:30:00      562   2026-10-25 02:39:22
IndexOptimize - USER_DATABASES - Daily     2026-10-25 02:30:00  -343599444  2015-12-05 06:12:36  <-- corrupto
IndexOptimize - USER_DATABASES - Weekend   2026-10-24 02:30:00    50360   2026-10-24 16:29:20  <--
IndexOptimize - USER_DATABASES - Daily     2026-10-23 02:30:08     3606   2026-10-23 03:30:14
IndexOptimize - USER_DATABASES - Daily     2026-10-22 02:30:00     5102   2026-10-22 03:55:02
DBA_CheckDB                                2026-10-25 18:00:02     1535   2026-10-25 18:25:37
\end{lstlisting}

\subsection{Recursos, hipervisor y alta disponibilidad}

\begin{lstlisting}[style=salida]
SELECT ram_mb, libre_mb, system_memory_state_desc FROM sys.dm_os_sys_memory
Time: 31 ms

ram_mb    libre_mb   system_memory_state_desc
131071    25657      Available physical memory is high

SELECT en_uso_mb, process_physical_memory_low ... FROM sys.dm_os_process_memory
Time: 16 ms

en_uso_mb   process_physical_memory_low   process_virtual_memory_low
97315       false                         false
\end{lstlisting}

\begin{lstlisting}[style=salida]
SELECT SUM(runnable_tasks_count), SUM(work_queue_count) ... dm_os_schedulers
Time: 0 ms

tareas_en_cola  sin_hilo  tareas_totales  workers  load_factor_max  load_factor_medio
5               0         151             146      19               13

SELECT TOP 10 ... FROM sys.dm_os_ring_buffers (RING_BUFFER_SCHEDULER_MONITOR)
Time: 62 ms

hora                      cpu_sql   cpu_otros   cpu_libre
2026-10-29 11:35:00.503        65           4          31
2026-10-29 11:36:00.697        75           6          19
2026-10-29 11:37:00.877        83           5          12
2026-10-29 11:38:01.063        77           4          19
2026-10-29 11:39:01.253        67           3          30
2026-10-29 11:40:01.440        59           3          38
2026-10-29 11:41:01.627        80           3          17
2026-10-29 11:42:01.813        79           3          18
2026-10-29 11:43:02.003        76           3          21
2026-10-29 11:44:02.190        61           3          36
\end{lstlisting}

\begin{lstlisting}[style=salida]
SELECT TOP 8 ... FROM sys.dm_io_virtual_file_stats(NULL, NULL) ...
Time: 47 ms

bd              fichero            tipo   lecturas     ms/lec   escrituras     ms/esc
ContaAlaplana   IN3ContaTpl_Log    LOG       65822          9      1063679           0
ContaSTN        IN3ContaTpl_Log    LOG      202724          7      3386531           0
ContaVitaCer    IN3Co...           LOG         ...          8          ...           1
D365Eventos     D365Eventos_log    LOG         ...          9          ...           1
Informatica     Informatica_Log    LOG         ...          8          ...          12
OMS             OMS_log            LOG         ...          8          ...           1
Grupo           In3Tplnew_log      LOG     9557386          7    281926025           2
IFR             IFR_log            LOG    22022516          7    363247671           2
\end{lstlisting}

\begin{lstlisting}[style=salida]
SELECT ... FROM sys.dm_db_log_info(DB_ID(...))  /  tempdb.sys.dm_db_file_space_usage
Time: 31 ms + 16 ms

bd        vlfs   activos   tam_medio_mb   log_total_gb
Cenusa      97        82          110.0           10.4
Grupo     1940        12           63.9          121.0
IFR        127         6         1465.7          181.8

version_store_mb   objetos_usuario_mb   objetos_internos_mb   libre_mb
31                 181                  1268                  430726
\end{lstlisting}

\begin{lstlisting}[style=salida]
GET /iwc/rest/databases   (vinculacion con vCenter)

id  name                          linkedToVm  vm             host ESXi     vcsa
3   AQUILES                       true        Aquiles        10.150.0.2    vcsa.central.local
2   PROMETEO                      true        Prometeo       10.150.0.9    vcsa.central.local
5   SRVPBI                        true        SRVPBI         10.150.0.9    vcsa.central.local
1   SRVPRORPDB01 via ATENEA2      true        SRVPRORPDB01   10.150.0.8    vcsa.central.local
4   SRVPROSCDB1 via SCALEPRODBAO  true        SRVPROSCDB1    10.150.0.27   vcsa.central.local

GET /iwc/rest/databases/1/metrics/data?metrics=...  (granularidad 600 s, 7 puntos)

metrica                   serie           media       maximo
vmReadyTime               SRVPRORPDB01      0,094       0,13
vmCpuCoStopPercent        SRVPRORPDB01      0           0
vmPercentBallooned        SRVPRORPDB01      0           0
vmBalloonSize             SRVPRORPDB01      0           0
vmMemorySwapInRate        SRVPRORPDB01      0           0
vmMemorySwapOutRate       SRVPRORPDB01      0           0
vmMemoryGranted           SRVPRORPDB01    130778      130778
vmMemoryActive            SRVPRORPDB01     71,16       74,47
vmCpuUsagePercent         SRVPRORPDB01     96,38      114,80
hostCpuUsagePercent       10.150.0.8       40,16       45,52
hostMemoryUsage           10.150.0.8       39,35       39,37
hostMaxTotalDiskLatency   10.150.0.8       59         211        <- artefacto

Desglose por dispositivo del host 10.150.0.8 -- descarta el artefacto anterior

dispositivo (DGC Fibre Channel Disk)        ms lectura   ms escritura
naa.6006016051bc51001367656699695a40            0,0          0,0
naa.6006016051bc51005f6d65661aaa6803            0,0          0,0
naa.6006016051bc51002c6b6566ef2fd7fe            0,1          0,9
naa.6006016051bc51008f6f6566a55ec291            0,5          0,0
naa.6006016051bc51008ec6916839933029            0,6          0,1
naa.6006016051bc5100086f65669cc587fa            0,8          0,7
naa.6006016051bc5100a1bf9168335813e0            0,8          1,3
\end{lstlisting}

\begin{lstlisting}[style=salida]
SELECT ag.name, ar.replica_server_name, ar.availability_mode_desc ...
Time: 1844 ms

grupo        replica        modo_commit            failover  rol        salud
AG-ATENEA2   SRVPRORPDB01   ASYNCHRONOUS_COMMIT    MANUAL    PRIMARY    HEALTHY
AG-ATENEA2   SRVPRORPDB02   ASYNCHRONOUS_COMMIT    MANUAL    SECONDARY  HEALTHY

SELECT DB_NAME(drs.database_id), drs.synchronization_state_desc ...
Time: 312 ms

bd         replica        estado         cola_envio  cola_redo  tasa_redo  retraso_s
ContaSTN   SRVPRORPDB02   SYNCHRONIZING          60          0      33201          0
Grupo      SRVPRORPDB02   SYNCHRONIZING          60         31      65609          0
IFR        SRVPRORPDB02   SYNCHRONIZING          60       2315      49590          0
OMS        SRVPRORPDB02   SYNCHRONIZING           0         69      33049          0
\end{lstlisting}

% =====================================================================
\newpage
\section{Anexo: planificaciones de menos de una hora}
\label{sec:anexo}

Hoja de trabajo para el escalonado (acción 7). Solo las planificaciones
activas con repetición en segundos o minutos; quedan fuera las horarias y las
de una vez al día.

\begin{longtable}{p{8.3cm}rp{1.2cm}p{1.5cm}p{1.5cm}}
\toprule
\textbf{Job} & \textbf{Cada} & \textbf{Unid.} & \textbf{Desde} & \textbf{Hasta} \\
\midrule
\endhead
stnSCALE\_720\_DATOS\_PESADA & 10 & seg & 00:00 & 23:59 \\
stnSCALE\_720\_DATOS\_PESADA\_ACTUALIZA\_MUELLE & 10 & seg & 05:45 & 22:00 \\
stnSCALE\_COPIA\_SHIPPING\_LOAD & 10 & seg & 00:00 & 23:59 \\
stnSCALE\_720\_DATOS\_PESADA\_ACTUALIZA\_TELEFONO & 30 & seg & 05:50 & 22:00 \\
stnSCALE\_INT\_ENT\_expediciones & 30 & seg & 05:30 & 22:30 \\
stnSCALE\_INT\_ENT\_expediciones & 30 & seg & 22:31 & 05:29 \\
stnSCALE\_INT\_ENT\_PRODUCCION & 30 & seg & 05:30 & 22:30 \\
stnSCALE\_INT\_ENT\_PRODUCCION & 30 & seg & 22:31 & 05:29 \\
stnSCALE\_INT\_ENT\_ORDENCARGA & 75 & seg & 05:30 & 22:00 \\
stnSCALE\_INT\_ENT & 90 & seg & 05:30 & 22:30 \\
stnSCALE\_INT\_ENT\_AJUSTES & 90 & seg & 05:30 & 22:30 \\
stnSCALE\_INT\_ENT\_ORDENCARGA & 100 & seg & 22:01 & 05:29 \\
stnSCALE\_INT\_ENT\_PARRILLA & 100 & seg & 00:00 & 23:59 \\
\midrule
D365\_ENVIO\_RECEPCIONES & 1 & min & 05:00 & 23:59 \\
STNAUTO\_AUTOMATIZAR\_ALBARANEADO & 1 & min & 00:00 & 23:59 \\
STNSCALE\_ERRORES\_ILSPRO & 1 & min & 00:00 & 23:59 \\
stnSCALE\_INT\_SAL & 1 & min & 05:30 & 22:30 \\
stnSCALE\_INT\_SAL\_ORDENCARGA & 1 & min & 05:30 & 22:30 \\
stnSCALE\_INT\_SAL\_PRODUCCION & 1 & min & 05:30 & 22:30 \\
D365\_SINCRONIZAR\_OMS\_DISP\_MAX & 2 & min & 05:00 & 23:59 \\
Interface APP WEB & 2 & min & 00:00 & 23:59 \\
RECALCULAR\_ALBARANES & 2 & min & 00:00 & 23:59 \\
REENVIAR ARTICULOS & 2 & min & 00:00 & 23:59 \\
REVISAR BLOQUEOS DE 7:30 a 8:30 & 2 & min & 06:00 & 09:00 \\
stn\_regenera\_op\_modificadas & 2 & min & 00:00 & 23:59 \\
stnSCALE\_EnvioTransportistas & 2 & min & 00:00 & 23:59 \\
ASPROVA\_COPIAR\_TRABAJO\_REAL & 3 & min & 00:00 & 23:59 \\
D365\_SYNC\_CREAR\_PEDIDO\_BLOQUEO\_DESDE365 & 3 & min & 00:00 & 23:59 \\
D365\_SYNC\_TRANSFORMAR\_DISP\_MAX & 3 & min & 05:00 & 23:59 \\
PREVISION\_CARGA\_CARGADATOS & 3 & min & 06:00 & 22:00 \\
STN\_OC\_ENPROCESO\_CARGA\_CARGADATOS & 3 & min & 00:00 & 23:59 \\
STNSCALE\_ERRORES & 3 & min & 00:00 & 23:59 \\
stnSCALE\_INT\_ENT & 3 & min & 22:31 & 05:29 \\
stnSCALE\_INT\_ENT\_AJUSTES & 3 & min & 22:31 & 05:29 \\
stnSCALE\_INT\_SAL & 3 & min & 22:31 & 05:29 \\
stnSCALE\_INT\_SAL\_PRODUCCION & 3 & min & 22:31 & 05:29 \\
D365\_SYNC\_CREAR\_OC\_GRUPO\_DESDE\_D365 & 4 & min & 00:00 & 23:59 \\
D365\_SYNC\_CREAR\_PEDIDO\_GRUPO\_DESDED365 & 4 & min & 00:00 & 23:59 \\
D365\_SYNC\_MANTENER\_RESERVAS\_GRUPO\_DESDED365 & 4 & min & 00:00 & 23:59 \\
stnscale\_albaranear\_mensajeria & 4 & min & 18:00 & 00:15 \\
D365\_SINCRONIZAR\_DATOS\_COMPRAS & 5 & min & 07:00 & 22:00 \\
D365\_SYNC\_MANTENER\_PEDIDO\_GRUPO\_DESDED365 & 5 & min & 00:00 & 23:59 \\
DESACCIONAR ORDENES & 5 & min & 05:00 & 23:00 \\
RECALCULAR\_PEDIDOS & 5 & min & 08:00 & 20:59 \\
REVISAR BLOQUEOS DE 7:30 a 8:30 & 5 & min & 16:55 & 17:40 \\
REVISAR BLOQUEOS DE 7:30 a 8:30 & 5 & min & 01:55 & 02:20 \\
STN\_ERP\_StnCRMSOLICITUD\_AprobacionAuto & 5 & min & 08:00 & 20:00 \\
STN\_SINCRONIZASTOCKRUMBO & 5 & min & 06:00 & 22:45 \\
stnSCALE\_720\_DATOS\_PESADA\_ACTUALIZA\_MUELLE & 5 & min & 22:00 & 05:44 \\
stnSCALE\_720\_DATOS\_PESADA\_ACTUALIZA\_TELEFONO & 5 & min & 22:00 & 05:49 \\
stnSCALE\_720\_DATOS\_PESADA\_BORRADO & 5 & min & 00:00 & 23:59 \\
stnSCALE\_EnvioMosaicoCarga & 5 & min & 00:00 & 23:59 \\
stnSCALE\_INT\_SAL\_ORDENCARGA & 5 & min & 22:31 & 05:29 \\
ACTUALIZA\_OPERACION & 10 & min & 00:00 & 23:59 \\
CENUSA - ARREGLA SERVICO OC CENUSA & 10 & min & 00:00 & 23:59 \\
Copiar SN\_PARRILLA\_LOG a Grupo & 10 & min & 00:00 & 23:59 \\
D365\_SYNC\_DATOS\_CLIENTE & 10 & min & 00:00 & 23:59 \\
STN\_MARCAR\_LINEAS\_PREPARADO\_TRASLADADO & 10 & min & 00:00 & 23:59 \\
stnscale\_albaranear\_mensajeria & 10 & min & 09:00 & 22:00 \\
ACT Estadisticas carga LEANM & 15 & min & 06:00 & 23:59 \\
ACTUALIZAR CODIGO UNICO & 15 & min & 08:00 & 16:45 \\
ASIGNACION AUTOMATICA CROSS-PARAGUAS & 15 & min & 06:00 & 21:00 \\
COMPROBAR AJUSTES STOCK & 15 & min & 00:00 & 23:59 \\
CONTROL JOBS SINFIN & 15 & min & 00:00 & 23:59 \\
D365\_SINCRONIZAR\_OMS\_DISP\_NORM & 15 & min & 07:00 & 23:59 \\
PREVISION\_CARGA\_CARGADATOS & 15 & min & 22:00 & 05:59 \\
STN\_ERP\_CARGAR\_ANALISIS\_VENTAS\_DIARIO & 15 & min & 06:00 & 23:00 \\
stnSCALE\_Revisar\_pallets\_SCALE & 15 & min & 00:00 & 23:59 \\
D365\_SYNC\_TRANSFORMAR\_DISP\_NORM & 18 & min & 05:00 & 23:59 \\
ACTUALIZAR\_FECHACOBRO\_EFECTOS\_DESDE\_CONCILIADOS & 30 & min & 08:00 & 20:00 \\
DesbloquearEmpresasSaldoNegativo & 30 & min & 08:00 & 19:00 \\
IN3 RUMBO ENVIO & 30 & min & 09:00 & 18:00 \\
stn720\_LIMPIA\_DATOS\_PESADA & 30 & min & 00:00 & 23:59 \\
stn\_crear\_stn\_bloqueo\_material\_producido & 30 & min & 00:00 & 23:59 \\
stn\_pzpedirproveedor\_CARGAR\_DATOS & 30 & min & 08:00 & 19:00 \\
STNSCALE\_ANULACIONES\_INCOMPLETAS & 30 & min & 00:00 & 23:59 \\
stnSCALE\_EnviaDatosEntregaPedido & 30 & min & 00:00 & 23:59 \\
stnSCALE\_EnviaDatosFechaEntregaPedido & 30 & min & 00:00 & 23:59 \\
stnSCALE\_EnviaDatosSuPedidoWickes & 30 & min & 00:00 & 23:59 \\
D365\_RECUPERAR\_PERDIDOS\_DATAVERSE\_CLAUDIA & 60 & min & 08:30 & 23:30 \\
\bottomrule
\end{longtable}

\noindent \textbf{80 planificaciones de menos de una hora.} A estas hay que
sumar 56 horarias y más de 40 diarias, hasta las 140 sub-diarias citadas en
el informe.

% =====================================================================
\newpage
\section{Glosario}

\begin{longtable}{p{4.6cm}p{11.2cm}}
\toprule
\textbf{Término} & \textbf{Qué significa} \\
\midrule
\endhead
CXCONSUMER, CXPACKET, CXSYNC\_PORT, EXECSYNC
  & Esperas de \emph{paralelismo}: un hilo espera a que otros hilos de la
    misma consulta terminen su parte. Mucho tiempo aquí indica que se está
    repartiendo trabajo que no compensa repartir \\
SOS\_SCHEDULER\_YIELD
  & La sesión agotó su turno de CPU y cede voluntariamente. Indica
    competencia por procesador, no por bloqueos \\
LCK\_M\_*
  & Espera real por un bloqueo de objeto. \cod{LCK\_M\_SCH\_M} es bloqueo de
    modificación de esquema, el que toman \cod{ALTER TABLE},
    reconstrucciones de índice y \cod{TRUNCATE} \\
PAGEIOLATCH\_SH
  & Espera a que una página se lea de disco a memoria. Sube cuando falta
    memoria para el \emph{buffer pool} \\
Espera inactiva (\emph{benign wait})
  & Espera de un hilo de fondo que no tiene trabajo. \cod{HADR\_WORK\_QUEUE},
    \cod{HADR\_TIMER\_TASK} o \cod{SLEEP\_*}. Deben excluirse del análisis:
    acumulan cifras enormes sin significar nada \\
MAXDOP
  & \emph{Max degree of parallelism}. Número máximo de núcleos que puede usar
    una sola consulta. A nivel de base de datos, el valor 0 significa
    «heredar el de la instancia» \\
\emph{Cost threshold for parallelism}
  & Coste estimado a partir del cual el optimizador considera un plan
    paralelo. Opera sobre el coste \emph{estimado}: si la estimación es mala,
    el umbral no protege \\
Concesión de memoria (\emph{memory grant})
  & Memoria que el optimizador reserva para una consulta antes de
    ejecutarla, calculada sobre las filas que \emph{cree} que manejará. Si se
    concede mucho y se usa poco, la estimación era errónea \\
\emph{Page life expectancy} (PLE)
  & Segundos que una página permanece en memoria antes de ser desalojada.
    Valores bajos indican presión de memoria \\
Escalado de bloqueos
  & Conversión automática de muchos bloqueos de fila en un único bloqueo de
    tabla, a partir de unas 5.000 filas. Ahorra memoria pero convierte una
    escritura pequeña en un bloqueo total de la tabla \\
RCSI (\cod{READ\_COMMITTED\_SNAPSHOT})
  & Opción de base de datos que hace que los lectores usen versiones de fila
    guardadas en TempDB en lugar de tomar bloqueos compartidos. Elimina la
    mayor parte del bloqueo lector-escritor \\
\emph{Version store}
  & Zona de TempDB donde se guardan las versiones de fila que usa RCSI \\
Disparador (\emph{trigger})
  & Código que se ejecuta automáticamente ante un \cod{INSERT},
    \cod{UPDATE} o \cod{DELETE}. \textbf{Corre dentro de la transacción que
    lo dispara}, de modo que alarga el tiempo que esa transacción retiene sus
    bloqueos \\
\emph{Ready time}
  & Porcentaje de tiempo que una máquina virtual está lista para ejecutar
    pero espera a que el hipervisor le asigne CPU física. Por encima del
    5\,\% indica sobresuscripción del host \\
\emph{Co-stop}
  & Tiempo que unos vCPU de una VM esperan a otros para poder avanzar
    sincronizados. Relevante en máquinas con muchos vCPU \\
\emph{Ballooning}
  & Mecanismo por el que el hipervisor reclama memoria de una VM cuando el
    host va justo. Si es mayor que cero, la memoria «libre» del invitado no
    es real \\
VLF (\emph{Virtual Log File})
  & Fragmento interno del registro de transacciones. Demasiados ralentizan
    arranque y recuperación; demasiado grandes impiden truncar el log con
    precisión \\
\bottomrule
\end{longtable}

\vfill
\begin{center}
\color{gris}\footnotesize
Informe generado el 29 de octubre de 2026 · Versión 4.0 consolidada\\
Servicios Generales de la Plana · Departamento de IT
\end{center}

\end{document}
